\documentclass[11pt]{amsart}
\usepackage[T1]{fontenc}
\usepackage{lmodern,amsmath,amssymb,mathtools}
\usepackage{tikz-cd,microtype}
\usepackage[a4paper,margin=30mm]{geometry}
\usepackage[colorlinks=true,linkcolor=blue!50!black,citecolor=blue!50!black,urlcolor=blue!50!black]{hyperref}
\def\UrlBreaks{\do\/\do-\do.}
\numberwithin{equation}{section}
\newtheorem{theorem}{Theorem}[section]
\newtheorem{lemma}[theorem]{Lemma}
\newtheorem{proposition}[theorem]{Proposition}
\newtheorem{corollary}[theorem]{Corollary}
\theoremstyle{definition}
\newtheorem{definition}[theorem]{Definition}
\newtheorem{remark}[theorem]{Remark}
\newtheorem{example}[theorem]{Example}
\newcommand{\Q}{\mathbb Q}
\newcommand{\R}{\mathbb R}
\newcommand{\C}{\mathbb C}
\newcommand{\Z}{\mathbb Z}
\newcommand{\OO}{\mathcal O}
\newcommand{\PP}{\mathbb P}
\DeclareMathOperator{\Pic}{Pic}
\DeclareMathOperator{\Supp}{Supp}
\DeclareMathOperator{\rk}{rk}
\DeclareMathOperator{\ord}{ord}
\DeclareMathOperator{\mult}{mult}
\DeclareMathOperator{\vol}{vol}
\DeclareMathOperator{\Sym}{Sym}
\DeclareMathOperator{\Bir}{Bir}
\DeclareMathOperator{\codim}{codim}
\title[Numerical semiampleness of nef adjoint divisors]{Numerical Semiampleness of Nef Adjoint Divisors}
\author{OpenAI}
\date{October 3, 2026}
\subjclass[2020]{Primary 14E30; Secondary 14C20, 14J32}
\keywords{generalised abundance, numerical semiampleness, nef divisor, minimal model, Frobenius morphism}
\hypersetup{pdftitle={Numerical Semiampleness of Nef Adjoint Divisors},pdfauthor={OpenAI}}
\begin{document}
\begin{abstract}
We prove the Generalised Abundance Conjecture: if $(X,B)$ is a projective
klt $\Q$-pair over an algebraically closed field of characteristic zero,
$K_X+B$ is pseudo-effective, $M$ is a nef $\Q$-Cartier divisor on $X$,
and $K_X+B+M$ is nef, then $K_X+B+M$ is numerically equivalent to a
semiample $\Q$-Cartier divisor on $X$.
\end{abstract}
\maketitle
\setcounter{tocdepth}{1}
\tableofcontents
\section{Introduction}
\label{sec:introduction}

A nef divisor has nonnegative degree on every curve. A semiample
divisor has a positive Cartier multiple generated by global sections,
and therefore defines a morphism to a projective variety. The passage
from numerical positivity to such a morphism is central to the minimal
model program. For an ordinary log canonical divisor, this is the
abundance problem. Adding an arbitrary nef divisor changes the natural
conclusion: one must allow replacement by a numerically equivalent
divisor.

We call a rational Cartier divisor \emph{numerically semiample} if it is
numerically equivalent to a semiample rational Cartier divisor on the
same variety. Here numerical equivalence means equality of degrees on
all integral curves. We prove the following statement.

\begin{theorem}\label{thm:main}
Let $(X,B)$ be a projective klt $\Q$-pair over an algebraically closed
field of characteristic zero. Suppose that $K_X+B$ is pseudo-effective,
and let $M$ be a nef $\Q$-Cartier divisor on $X$. If
\[
                         D=K_X+B+M
\]
is nef, then there is a semiample $\Q$-Cartier divisor $L$ on $X$ such
that $D\equiv L$.
\end{theorem}

The rationality assumptions apply to the boundary and the nef summand;
in particular the latter is a divisor on $X$ itself. The proof uses the
log-abundance theorem of \cite[Theorem~1.1]{LA} and the terminating
minimal-model programs of \cite[Theorem~1.1 and Section~3.3]{MM}.
Their precise forms are stated in Section~\ref{sec:preliminaries}.

Theorem~\ref{thm:main} contains the Generalised Abundance Conjecture
formulated by Lazi\'c--Peternell \cite{LP}, where the nef summand is
Cartier. Numerical equivalence is necessary even in dimension one.
Indeed, a nontorsion degree-zero line bundle on a complex elliptic curve
is nef and numerically trivial, but none of its positive powers has a
nonzero section. Since the canonical bundle of the curve is trivial,
this is an example with $B=0$ in Theorem~\ref{thm:main}.
For the ordinary adjoint, Fukuda's theorem
\cite[Theorem~0.1]{Fukuda} shows that a semiample numerical
representative makes $K_X+B$ itself semiample. For $M=0$, the conclusion
is therefore equivalent to ordinary klt abundance, which is an input
to this proof.

\begin{corollary}\label{cor:cy}
Let $(X,B)$ be a projective klt $\Q$-pair in characteristic zero with
$K_X+B\equiv0$. Every nef $\Q$-Cartier divisor on $X$ is numerically
semiample.
\end{corollary}
\begin{proof}
For a nef divisor $M$, the divisor $K_X+B+M$ is nef and numerically
equivalent to $M$. Apply Theorem~\ref{thm:main}.
\end{proof}

\subsection{Background and the role of the present argument}

The generalized-pair formalism retains a nef divisor on a higher
birational model as part of the adjoint data. Birkar--Zhang \cite{BZ}
developed effective birationality for these polarized pairs, extending
the tools used to study Iitaka fibrations. The canonical bundle formula
of Ambro \cite{Ambro} is one source of such nef data. In this paper the
original nef summand descends to $X$, but a minimal-model program forces
us to keep track of it as a fixed nef divisor on a higher model.

Lazi\'c--Peternell \cite{LP} formulated generalized abundance as a
numerical extension of abundance and related it to semiampleness on
varieties with trivial canonical class. Their Theorem~B, assuming
termination of flips and abundance in dimensions at most $n$, proves
numerical semiampleness when the ordinary adjoint has positive numerical
dimension. Its numerical-dimension-zero case also assumes their
Semiampleness Conjecture. Their Corollaries~C and~D give unconditional
results for surfaces and for threefolds with ordinary adjoint of positive
numerical dimension. These results identify the main remaining
positivity issue: an arbitrary nef summand on a variety with trivial
canonical class need not supply sections directly.

Subsequent work has addressed numerical nonvanishing, which asks for
an effective numerical representative, and refinements of numerical
semiampleness. Lazi\'c--Peternell \cite[Theorems~B and~C]{LPII}
obtained numerical nonvanishing under additional numerical or metric
hypotheses and, in the metric case, conjectural MMP and abundance
assumptions. For a Cartier nef summand, Chaudhuri
\cite[Corollary~3]{Chaudhuri} proved generalized abundance when the
nef-reduction base of $D$ has dimension two and $K_X+B$ is effective,
or when that base has dimension three and $\kappa(K_X+B)>0$.
Here nef reduction contracts the curves of degree zero through very
general points; its precise properties are recalled in
Section~\ref{sec:preliminaries}. On surfaces, Fontanari
\cite[Theorem~1.4]{Fontanari} proved actual semiampleness when the nef
summand is Cartier and its numerical class is not proportional to
that of the ordinary adjoint.

The birational reductions below follow \cite[Sections~4--7]{LP}.
The companion results \cite{LA,MM} supply ordinary log abundance and
the existence of the particular terminating programs required in those
reductions. The further ingredient proved here is that a nef divisor
of full nef dimension on a klt Calabi--Yau pair is big. Full nef
dimension means that every curve through a very general point has
positive degree. Bigness is stronger: it requires the volume of the
divisor to be positive. The gap between these two conditions is the
geometric issue addressed by the proof.

Our treatment adapts the moving-jet and Frobenius arguments of
\cite[Sections~7--9]{LA}. The moving-point method has antecedents in
the local-positivity work of Ein--K\"uchle--Lazarsfeld \cite{EKL}.
The finite-data comparison uses Frobenius jets, as studied by
Musta\c t\u a--Schwede \cite{MustataSchwede}, the canonical filtrations
of Sun and Kitadai--Sumihiro \cite{Sun,KS}, and Langer's control of
Frobenius slopes \cite{Langer}. We reproduce the comparison proof
with its uniformity conditions. The adaptation to an arbitrary nef
summand occurs in the geometric exclusions preceding it: the polarized
birationality theorem of Birkar--Zhang gives curves of bounded degree,
contradicting a growing integral lower bound for their degrees. We
give explicit proofs of those exclusions, including the case of
correspondences between finite covers.

\subsection{Structure of the proof}

There are two simultaneous induction statements. The first says that
$K_X+B+cM$ is big for every rational $c>0$ whenever one positive
combination of $K_X+B$ and $M$ has positive degree on all curves through
a very general point. The second says that the positive part of
$K_X+B+M$ on a suitable smooth birational model is rational and
numerically semiample. The positive part is the divisor left after
removing the asymptotically fixed prime divisors. When the original
adjoint is nef, it has no such negative part, and the second assertion
descends to Theorem~\ref{thm:main}.

Section~\ref{sec:reduction} sets up this induction. A pair of compatible
minimal-model programs retains both a semiample ordinary adjoint and a
nearby nef generalized adjoint. Positive Iitaka dimension then permits
induction on the fibres. The remaining full-dimension case is a klt
Calabi--Yau pair.

Sections~\ref{sec:jet-tools} and~\ref{sec:geometry} treat the essential
terminal case with trivial canonical bundle. If a nef Cartier divisor
$L$ of full nef dimension were not big, its ample approximations could
be rescaled to bounded volume while their degrees on very general
curves tend to infinity. The bounded-degree obstruction controls the
base loci of moving jet kernels. It produces a small vanishing-order
bound on the variety itself and a large jet system on a projective
bundle over its square. Section~\ref{sec:frobenius} shows that these two
conclusions are incompatible. The comparison is made on fixed complex
data before reduction to positive characteristic.

Section~\ref{sec:completion} passes through an index-one cover and a
boundary perturbation to obtain the full-dimension assertion for all
klt Calabi--Yau pairs. Section~\ref{sec:positive} proves the
positive-part assertion by nef reduction and the canonical bundle
formula, then descends numerical semiampleness to the original variety
and to any algebraically closed field of characteristic zero.

\section{Conventions and birational inputs}
\label{sec:preliminaries}

Varieties are integral and projective over an algebraically closed
field of characteristic zero unless stated otherwise. A $\Q$-pair
$(X,B)$ consists of a normal variety, an effective rational boundary
$B$, and a rational Cartier divisor $K_X+B$. We use log discrepancies:
on a resolution $h:Y\to X$, the equality
$K_Y+B_Y=h^*(K_X+B)$ assigns log discrepancy $1-b_E$ to a prime divisor
with coefficient $b_E$ in $B_Y$. The pair is klt if all log
discrepancies are positive, and log canonical if they are nonnegative.
We use the standard cone, contraction, negativity, and resolution
theorems as in \cite{KM,BCHM}.

The symbols $\sim_\Q$ and $\equiv$ denote rational linear and numerical
equivalence. A divisor is pseudo-effective if its numerical class is
in the closure of the cone generated by effective divisors, and big if
its class lies in the interior of that cone. Its Iitaka dimension is
denoted by $\kappa$. Over $\C$, a very general point means a point
outside a specified countable union of proper closed subsets.

For a nef divisor $N$, its nef reduction is an almost holomorphic
dominant rational map with connected compact general fibres, on which
$N$ is numerically trivial, such that every curve through a very
general point not contracted by the map has positive $N$-degree
\cite{NefReduction}. Its base dimension is the \emph{nef dimension}
$n(X,N)$. Thus $n(X,N)=\dim X$ precisely when every curve through a
very general point has positive $N$-degree. This numerical condition
does not assert bigness.

For a pseudo-effective real divisor $D$ on a smooth variety and a
prime divisor $\Gamma$, fix an ample divisor $A$ and set
\[
 \sigma_\Gamma(D)=\lim_{\delta\downarrow0}
 \inf\{\mult_\Gamma D':D'\ge0,\ D'\sim_\R D+\delta A\}.
\]
These numbers are independent of $A$ and depend only on the numerical
class of $D$. The divisorial Zariski decomposition is
\[
 N_\sigma(D)=\sum_\Gamma\sigma_\Gamma(D)\Gamma,
 \qquad P_\sigma(D)=D-N_\sigma(D).
\]
The sum has finite support. A nef divisor has $N_\sigma(D)=0$.
Further properties from \cite{Nakayama,LP} are recorded at their point
of use in Section~\ref{sec:positive}.

\begin{lemma}\label{lemma:pseudo-vg}
Let $D$ be a pseudo-effective rational Cartier divisor on a projective
complex variety $X$. Every integral curve through a sufficiently very
general point has nonnegative $D$-degree.
\end{lemma}
\begin{proof}
Fix an ample Cartier divisor $A$ and positive rational numbers
$\delta_i\to0$. Each $D+\delta_iA$ is big, so there is an effective
rational divisor $E_i\sim_\Q D+\delta_iA$. Choose a point outside
$\bigcup_i\Supp E_i$. A curve through it is contained in none of
these supports; hence $(D+\delta_iA)\cdot C\ge0$ for every $i$.
Let $i\to\infty$.
\end{proof}

We state explicitly the companion results used throughout. We require
existence of a terminating choice of program, rather than termination
of every sequence of flips.

\begin{theorem}[Log abundance, {\cite[Theorem~1.1]{LA}}]
\label{input:abundance}
Let $(X,B)$ be a projective log canonical $\Q$-pair over an
algebraically closed field of characteristic zero. If $K_X+B$ is nef,
then it is semiample.
\end{theorem}

For the next input a nef rational Cartier b-divisor is specified by a
nef rational Cartier divisor on a projective birational model; it is
pulled back unchanged on all higher models. Its trace on a lower model
is its pushforward and need not be nef. Generalized discrepancies are
defined by pulling back the adjoint while keeping this fixed nef data.

\begin{theorem}[Minimal-model programs,
{\cite[Theorem~1.1 and Section~3.3]{MM}}]
\label{input:models}
Let $(X,B+\mathbf M)$ be a projective generalized klt $\Q$-pair over $\C$,
with $X$ $\Q$-factorial and with fixed nef rational Cartier b-divisor
data. If $K_X+B+M_X$ is pseudo-effective, there is a terminating choice
of its minimal-model program, with $\Q$-factorial generalized klt
models and a nef endpoint. The program extracts no divisors and
generalized log discrepancies do not decrease. On a common resolution
of its initial and final models, the pullback of the initial adjoint
is the pullback of the final adjoint plus an effective divisor
exceptional over the final model.
\end{theorem}

The ordinary klt case follows by taking the nef data to be zero. In
that case Theorem~\ref{input:abundance} makes the endpoint a good
minimal model. The moving-center and Frobenius results used from
\cite{LA} are stated separately in
Sections~\ref{sec:jet-tools} and~\ref{sec:frobenius}; neither theorem
above asserts abundance for a nonzero generalized nef summand.

\section{Reduction of full nef dimension to Calabi--Yau pairs}
\label{sec:reduction}

We work over $\C$ until the final descent argument.  The reductions in
this section follow the strategy of Lazi\'c--Peternell
\cite[Sections 4 and 5]{LP}.  We use the terminating programs in
Theorem~\ref{input:models}, together with the abundance input
Theorem~\ref{input:abundance}, to make the precise reductions needed below.
The essential point is to retain control of two adjoints: an ordinary
semiample adjoint and a nearby generalized nef adjoint.

\subsection{The induction statements}

For a projective klt $\Q$-pair $(X,B)$ and a nef $\Q$-Cartier divisor $M$
on $X$, write
\[
                         T=K_X+B.
\]
Suppose that $T$ is pseudo-effective.  We will use the condition that,
for some $a\in\Q_{>0}$,
\begin{equation}\label{eq:full-condition}
 (T+aM)\cdot C>0
 \quad\text{for every integral curve $C$ through a very general point of $X$.}
\end{equation}
The divisor $T+aM$ need not be nef.  By
Lemma~\ref{lemma:pseudo-vg}, Condition~\eqref{eq:full-condition} holds
whenever $M$ has full nef dimension.  When $T+aM$ is nef, the condition
says exactly that its nef dimension is $\dim X$.

We prove the following two assertions together, by induction on $n$.
The notation $P_\sigma$ denotes the positive part of the divisorial
Zariski decomposition; its properties used in the proof are recalled
in Section~\ref{sec:positive}.
\begin{description}
\item[$(P_n)$] If $(X,B)$ is a projective klt $\Q$-pair of dimension $n$,
$T=K_X+B$ is pseudo-effective, $M$ is a nef $\Q$-Cartier divisor on $X$,
and \eqref{eq:full-condition} holds, then $T+cM$ is big for every
$c\in\Q_{>0}$.
\item[$(Z_n)$] If $(X,B)$ and $M$ have the same hypotheses, without
\eqref{eq:full-condition}, then there is a projective birational morphism
$w\colon W\to X$ from a smooth projective variety such that
\[
                         P_\sigma\bigl(w^*(T+M)\bigr)
\]
is a rational divisor numerically equivalent to a semiample
$\Q$-divisor on $W$.
\end{description}

Both assertions hold in dimensions zero and one.  In dimension one,
pseudo-effectivity of $T$ and nefness of $M$ mean that their degrees are
nonnegative.  Condition~\eqref{eq:full-condition} makes every positive
combination have positive degree, hence be ample.  A divisor of degree
zero is numerically equivalent to zero, and a divisor of positive degree
is ample.  These observations also give $(Z_1)$, since a nef divisor has
zero negative part.  On a point the assertions use the usual conventions.

Fix $n\geq2$, and assume $(P_j)$ and $(Z_j)$ for $j<n$.
At any stage we may replace $X$ by a small projective
$\Q$-factorialization.  The ordinary adjoint pulls back crepantly, $M$
pulls back to a nef divisor, and \eqref{eq:full-condition} persists by
using points in the isomorphism locus.  Bigness descends under this
birational pullback.  A smooth model proving $(Z_n)$ for the
$\Q$-factorialization also proves it for $X$.  We therefore impose
$\Q$-factoriality when running the programs below.

\subsection{Keeping a semiample adjoint fixed}

We first record the elementary ray comparison used in the second
program.  The nef part of a generalized pair in this argument is the
fixed nef $\Q$-Cartier b-divisor determined by $M$ on the original
variety.  Its trace on a later model need not be nef.

\begin{lemma}\label{lemma:generalized-ray-bound}
Let $(V,\Delta+\mathbf N)$ be a projective generalized klt $\Q$-pair
of dimension $n$, with effective boundary, $\Q$-factorial $V$, and nef
rational Cartier b-divisor data.  Every
$(K_V+\Delta+N_V)$-negative extremal ray has an integral rational curve
generator $C$ for which
\[
                -4n\leq (K_V+\Delta+N_V)\cdot C<0.
\]
The contraction, and the flip when it is small, can be realized using
an ordinary klt adjoint negative on the same ray.  Birational MMP
outputs remain $\Q$-factorial.  On a common resolution of a birational
step, the pullback of the input generalized adjoint equals the pullback
of its output transform plus an effective divisor exceptional over
the output.
\end{lemma}

\begin{proof}
Set $G=K_V+\Delta+N_V$, and fix a $G$-negative extremal ray.
For an ample rational divisor $A$ on $V$, choose $\eta>0$ rational and
small enough that $G+\eta A$ is still negative on the ray and
\[
                 -(G+\eta A)\cdot\gamma
                      \geq -\tfrac12G\cdot\gamma
\]
for its nonzero classes $\gamma$.  We may represent $G+\eta A$ as an
ordinary klt adjoint.  Here is the reason this remains valid when the
nef b-divisor does not descend to $V$.

Choose a projective log resolution $h\colon V'\to V$ carrying its nef
part $N'$ and write
\[
                  K_{V'}+\Delta'+N'=h^*G.
\]
The coefficients of $\Delta'$ are strictly less than one.  The
resolution can be chosen with an effective exceptional divisor $F$
such that $-F$ is $h$-ample.  For sufficiently small rational
$\varepsilon>0$, the divisor
$N'+\eta h^*A-\varepsilon F$ is ample.  A general effective rational
representative of this ample class can be chosen with sufficiently
small coefficients that, after adding $\varepsilon F$, it preserves
the sub-klt condition of $\Delta'$.  Its pushforward, added to
$\Delta$, is effective and defines a klt boundary $\Theta$ on $V$ with
\[
                       K_V+\Theta\sim_\Q G+\eta A.
\]
The crepant equality can be checked upstairs: its difference is
exceptional and relatively numerically trivial, hence zero by
negativity.

The ordinary cone theorem now gives a rational curve generator $C$
with $0<-(G+\eta A)\cdot C\leq2n$.  The preceding comparison yields
$-G\cdot C\leq4n$.  The contraction $c\colon V\to S$ is the
contraction of the same ray for this ordinary klt pair.  Since its
relative Picard number is one, there is a rational $\lambda>0$ such
that $G-\lambda(K_V+\Theta)$ is numerically trivial over $S$.
The Cartier descent statement for an ordinary klt extremal
contraction, after clearing denominators, gives a rational Cartier
divisor $D_S$ with
\[
                 G-\lambda(K_V+\Theta)\sim_\Q c^*D_S.
\]
In the small case, let $c^+\colon V^+\to S$ be the ordinary flip.
Taking strict transforms gives
\[
                 G^+\sim_\Q
                 \lambda(K_{V^+}+\Theta^+)+(c^+)^*D_S,
\]
which is ample over $S$.  Thus the ordinary flip is also the flip of
$G$.  The ordinary contraction and flip preserve $\Q$-factoriality
on the birational output.  Their effective exceptional adjoint
comparison, multiplied by $\lambda$, gives the asserted comparison
for $G$, since the terms pulled back from $S$ cancel.  With the nef
b-divisor fixed, this comparison is also the improvement of
generalized discrepancies used below.
\end{proof}

\begin{lemma}\label{lemma:two-nef-model}
Let $(X,B)$ be a projective $\Q$-factorial klt $\Q$-pair of dimension
$n$, let $T=K_X+B$ be pseudo-effective, and let $M$ be a nef
$\Q$-Cartier divisor on $X$.  Suppose
\eqref{eq:full-condition} holds.  Then there are a birational
contraction $X\dashrightarrow Y$, a number $s_0\in\Q_{>0}$, and a
common smooth projective resolution
\[
                       X\xleftarrow{p}W\xrightarrow{q}Y
\]
with the following properties.  If $B_Y$ and $N$ denote the transforms
of $B$ and $M$, and $U=K_Y+B_Y$, then
\begin{enumerate}
\item $(Y,B_Y)$ is klt and $\Q$-factorial, $U$ is semiample, and
$\kappa(Y,U)=\kappa(X,T)\geq0$;
\item $N$ is pseudo-effective, and $U+s_0N$ is nef of full nef
dimension;
\item there are effective $q$-exceptional rational divisors $E_0,E_1$
on $W$ such that
\begin{equation}\label{eq:two-comparisons}
 p^*T=q^*U+E_0,
 \qquad
 p^*(T+s_0M)=q^*(U+s_0N)+E_1.
\end{equation}
\end{enumerate}
\end{lemma}

\begin{proof}
By Theorem~\ref{input:models}, there is a terminating $T$-MMP, ending
on a $\Q$-factorial klt pair $(Y_0,B_0)$ with nef adjoint
$U_0=K_{Y_0}+B_0$.  Theorem~\ref{input:abundance} makes $U_0$
semiample.  The effective exceptional comparisons for the program
give
\[
                         \kappa(X,T)=\kappa(Y_0,U_0)\geq0.
\]
Let $N_0$ be the transform of $M$.  It is the trace on $Y_0$ of the
fixed nef b-divisor determined by $M$ on $X$.

Because this first program has finitely many steps, there is a
rational $s>0$ such that every one of them remains negative for the
transform of $T+uM$, for every rational $u\in(0,s]$.  Choose $s$ also
small enough to retain the strict sign on the flipped side of each
flip.  These are therefore programs for the corresponding generalized
klt pairs: they start generalized klt on $X$, and their discrepancies
improve at each step.  In particular, $(Y_0,B_0+s\mathbf M)$ is
generalized klt.  Pushforward to a $\Q$-factorial model preserves
pseudo-effectivity, so $N_0$ is pseudo-effective.

For completeness, this last numerical assertion applies also to
birational contractions with flips.  On a common resolution, numerical
pushforward to a $\Q$-factorial target is well defined: if a divisor
upstairs is numerically trivial, its difference from the pullback of
its pushforward is exceptional and relatively numerically trivial,
and the negativity lemma makes that difference zero.  Effective
approximation then proves preservation of pseudo-effectivity.

Choose $d>0$ so that $dU_0$ is Cartier and globally generated.  Choose
a rational $s_0\in(0,s)$ such that
\begin{equation}\label{eq:small-nef-coefficient}
                 \frac{1-s_0/s}{d}>4n\frac{s_0}{s}.
\end{equation}
The generalized adjoint $U_0+s_0N_0$ is pseudo-effective.  Apply
Theorem~\ref{input:models} to obtain a terminating program for it.
We claim that every step is $U_0$-trivial and is also negative for the
generalized adjoint with coefficient $s$, where throughout this
claim the symbols denote their transforms.

Suppose the claim holds up to the next step, and denote the current
transforms by $U_i,N_i$.  Then $dU_i$ is still Cartier and free, and
the pair with nef b-divisor coefficient $s$ is generalized klt.
For a $(U_i+s_0N_i)$-negative ray, the equality
\[
 U_i+s_0N_i=(1-s_0/s)U_i+(s_0/s)(U_i+sN_i)
\]
and nefness of $U_i$ show that the ray is $(U_i+sN_i)$-negative.
Choose its rational curve generator $C$ by
Lemma~\ref{lemma:generalized-ray-bound}.  If $U_i\cdot C>0$, then
integrality of $dU_i$ gives $U_i\cdot C\geq1/d$, whence
\[
 (U_i+s_0N_i)\cdot C
 \geq\frac{1-s_0/s}{d}-4n\frac{s_0}{s}>0,
\]
a contradiction.  Thus $U_i$ is trivial on the ray.

The morphism defined by $dU_i$ is constant on every fibre of the
extremal contraction and factors through it.  Its descended line
bundle is globally generated; pulling it to the flipped model, when
there is a flip, gives the transform of $dU_i$.  Hence that transform
is again Cartier and free, with equal pullbacks on a common
resolution.  Since the step is also negative for coefficient $s$,
generalized klt singularities with that coefficient persist.  This
proves the claim by induction through the finite program.

Let $Y$ be its output.  The ordinary adjoint $U$ is crepant throughout
the second program, so $(Y,B_Y)$ is klt, $U$ is semiample, and its
Iitaka dimension equals that of $T$.  The generalized adjoint
$U+s_0N$ is nef.  Combining the pullback comparisons of the two
programs gives \eqref{eq:two-comparisons} with effective
$q$-exceptional divisors.  The transform $N$ remains pseudo-effective.

It remains to prove full nef dimension.  Suppose instead that the
nef reduction of $U+s_0N$ has a positive-dimensional very general
compact fibre $F$.  Shrink the base of the almost holomorphic map so
that these compact fibres lie entirely in its morphism locus.  The
image in $Y$ of the exceptional locus of $q$ has
codimension at least two.  A general fibre $F$ meets each of its
components in codimension at least two, or not at all: discard the
images of components that do not dominate the nef-reduction base,
and use the fibre dimension theorem for the others.  We may
therefore choose a curve $C\subset F$ through a very general point,
avoiding this exceptional image.  If $\dim F>1$, take sufficiently
general complete intersections in $F$ through that point; if
$\dim F=1$, take $F$ itself.

Choose the point also in the common isomorphism locus of the
birational contraction and with the very general conditions of
\eqref{eq:full-condition} transported from $X$.  Since $U$ is nef
and $N$ is pseudo-effective, Lemma~\ref{lemma:pseudo-vg} gives
\[
                    U\cdot C=N\cdot C=0,
\]
because their positive combination has degree zero on $F$.
The strict lift $\widetilde C\subset W$ avoids $E_0$ and $E_1$.
Both equalities in \eqref{eq:two-comparisons} then give
\[
               p^*T\cdot\widetilde C
                =p^*M\cdot\widetilde C=0.
\]
The curve $p(\widetilde C)$ passes through the chosen very general
point of $X$ and contradicts \eqref{eq:full-condition}.
\end{proof}

\subsection{Positive Iitaka dimension and the remaining case}

The preceding construction leaves only one obstruction to bigness:
the ordinary semiample adjoint may define a map to a point.  If it
defines a positive-dimensional base, the induction hypothesis handles
its fibres.

\begin{proposition}\label{prop:full-positive}
Assume $(P_j)$ for $j<n$.  Let $(X,B)$ be a projective klt
$\Q$-pair of dimension $n$, let $T=K_X+B$ be pseudo-effective, and
let $M$ be a nef $\Q$-Cartier divisor on $X$.  If
\eqref{eq:full-condition} holds and $\kappa(X,T)>0$, then $T+cM$ is
big for every $c\in\Q_{>0}$.
\end{proposition}

\begin{proof}
After a small $\Q$-factorialization, apply
Lemma~\ref{lemma:two-nef-model}.  Let $h\colon Y\to S$ be the
contraction defined by the semiample divisor $U$.  There is an ample
rational divisor $A_S$ on $S$ with $U\sim_\Q h^*A_S$, and
$\dim S=\kappa(X,T)>0$.  Put $Q=U+s_0N$.

If a very general fibre $F$ has positive dimension, it has dimension
strictly less than $n$.  The restriction pair $(F,B_Y|_F)$ is klt
and
\[
                    K_F+B_Y|_F\sim_\Q U|_F\sim_\Q0.
\]
Moreover, $Q|_F$ is nef of full nef dimension.  Indeed, a very general
point of a very general $F$ can be chosen outside the countable union
excluded in the full-nef-dimension assertion for $Q$ on $Y$.
Since $U|_F\equiv0$, the divisor $N|_F$ is nef.  Apply $(P_{\dim F})$
to this restriction pair and $s_0N|_F$ to conclude that $Q|_F$ is big.
When the fibre dimension is zero, relative bigness is automatic.

Thus $Q$ is big over $S$.  To justify the passage to the generic
fibre, fix an ample Cartier divisor $H$ on $Y$ and shrink $S$ so
that $h$ is flat.  For every positive integer $\ell$ clearing the
denominators of $Q$, upper semicontinuity makes the locus
\[
        \bigl\{s\in S:
        H^0(F_s,\OO_{F_s}((\ell Q-H)|_{F_s}))\ne0\bigr\}
\]
closed.  Bigness on a very general fibre puts its base point in one
of these countably many loci.  Choose that point outside every
locus that is proper; the locus containing it must then be all of
$S$.  On the generic fibre, $\ell Q$ is consequently the sum of
the ample divisor $H$ and an effective divisor, which proves the
asserted relative bigness.  A relatively big divisor becomes big after adding
a sufficiently large pullback of an ample divisor on the base;
equivalently, apply the usual relative form of Kodaira's lemma
\cite[Lemma 3.2.1]{BCHM}.  It follows that
\[
                         Q+bU=(1+b)U+s_0N
\]
is big for some rational $b>0$.  Equation~\eqref{eq:two-comparisons}
then shows that $(1+b)T+s_0M$ is big.

Finally, fix $c>0$ rational and choose $\lambda>0$ rational with
$\lambda(1+b)<1$ and $\lambda s_0<c$.  Then
\[
 T+cM=\lambda\bigl((1+b)T+s_0M\bigr)
       +\bigl(1-\lambda(1+b)\bigr)T+(c-\lambda s_0)M.
\]
The first summand is big and the other two are pseudo-effective.
Their sum is big, as required.
\end{proof}

\begin{proposition}\label{prop:reduction-cy}
Assume $(P_j)$ for $j<n$.  Suppose that on every projective klt
$\Q$-pair $(V,\Delta)$ of dimension $n$ with
$K_V+\Delta\equiv0$, every nef $\Q$-Cartier divisor of full nef
dimension is big.  Then $(P_n)$ holds.
\end{proposition}

\begin{proof}
Use a small $\Q$-factorialization and
Lemma~\ref{lemma:two-nef-model}.  The nonnegative Iitaka dimension
of $T$ is either positive, in which case
Proposition~\ref{prop:full-positive} applies, or zero.  In the latter
case, semiampleness of $U$ implies $U\sim_\Q0$.  The pair $(Y,B_Y)$
is therefore a klt Calabi--Yau pair.  The nef divisor $U+s_0N$ has
full nef dimension, so it is big by the stated hypothesis.
The second equality in \eqref{eq:two-comparisons} gives bigness of
$T+s_0M$.  The positive-combination argument at the end of
Proposition~\ref{prop:full-positive} gives bigness of $T+cM$ for
every rational $c>0$.
\end{proof}

Consequently, the remaining work for $(P_n)$ is the full-nef-dimension
statement on klt Calabi--Yau pairs.  Theorem~\ref{thm:terminal-big}
establishes its terminal, zero-boundary form.  We then pass from that
form to all klt Calabi--Yau pairs in Section~\ref{sec:completion}.

\section{Bounded-degree curves and moving jet kernels}
\label{sec:jet-tools}

The geometric argument will produce subvarieties with bounded degree and
bounded canonical intersection.  We first explain how these bounds
contradict a growing lower bound for the degrees of curves.  We then
record the moving-center results from \cite{LA} that supply the required
subvarieties.  All varieties in this section are over $\C$.

\subsection{The bounded-degree obstruction}

\begin{lemma}[Bounded-degree curves]
\label{lemma:bounded-curves}
For every integer $f\geq1$ there is an integer $m_f>0$ with the
following property.  Let $Y$ be a smooth integral projective variety
of dimension $f$, let $D_0$ be a reduced simple normal crossing divisor
or the zero divisor, and let $L_0$ be a nef integral Cartier divisor
of full nef dimension on $Y$.  Suppose that $P_0$ is a nef rational
Cartier divisor and that $r,C_0,C_1$ are positive real numbers satisfying
\[
 0<P_0^f\leq C_0,\qquad P_0-rL_0\text{ is nef}.
\]
If $A_0=K_Y+D_0+L_0$ is big and
\begin{equation}
\label{j:adjoint-ratio}
 A_0\cdot P_0^{f-1}\leq C_1P_0^f,
\end{equation}
then
\begin{equation}
\label{j:bounded-curve-obstruction}
 r\leq(m_fC_1)^{f-1}C_0.
\end{equation}
The integer $m_f$ depends only on $f$, and not on the variety or
the divisors.  For dimensions $1\leq f\leq n$, one may replace all
$m_f$ by a single integer $m$ depending only on $n$.
\end{lemma}

\begin{proof}
The pair $(Y,D_0)$ is log canonical, its boundary coefficients belong
to $\{0,1\}$, and its nef polarizing divisor $L_0$ has Cartier index
one.  By \cite[Theorem~1.3]{BZ}, an integer $m_f$ depending only on
these data makes $|m_fA_0|$ birational.  Here $A_0$ is integral
Cartier, so no rounding or additional denominator is involved.

Resolve the base ideal by a projective birational morphism
$\mu:\widetilde Y\to Y$, with $\widetilde Y$ smooth, and write
\[
 \mu^*(m_fA_0)=H+F,
\]
where $H$ is the basepoint-free moving part and $F$ is effective.
The morphism defined by $H$ is birational onto its image; in
particular $H$ is nef and big.  In the rest of this proof, $P_0$
and $L_0$ denote their pullbacks to $\widetilde Y$.  Effectivity
of $F$ and nefness of $P_0$ give
\begin{equation}
\label{j:moving-adjoint-degree}
 H\cdot P_0^{f-1}\leq m_fC_1P_0^f.
\end{equation}

For $f\geq2$, put $s_i=P_0^{f-i}\cdot H^i$ for $0\leq i\leq f$.
Both divisors are nef and big, so the $s_i$ are positive.  The
Khovanskii--Teissier inequalities
\cite[Corollary~1.6.3(i) and Example~1.6.4]{LazarsfeldPAGI}
$s_i^2\geq s_{i-1}s_{i+1}$ show
that the ratios $s_i/s_{i-1}$ decrease.  Hence
\begin{equation}
\label{j:mixed-degree}
 P_0\cdot H^{f-1}
 \leq P_0^f\left(\frac{H\cdot P_0^{f-1}}{P_0^f}\right)^{f-1}
 \leq(m_fC_1)^{f-1}C_0.
\end{equation}
Choose a very general point $y\in\widetilde Y$ over the
isomorphism opens of both $\mu$ and the morphism defined by $H$.
We also choose $\mu(y)$ in the very general locus
that tests full nef dimension of $L_0$.  Pull back $f-1$ general
hyperplanes through the image of $y$ under the morphism defined
by $H$.  This subsystem has only the base point $y$, so the
general members meet properly and define an effective curve cycle
of class $H^{f-1}$; the component $C$ through $y$ is a curve
because their local equations are independent at $y$.
Nefness gives
\[
 P_0\cdot C\leq P_0\cdot H^{f-1}.
\]
The curve $C$ is not contracted by $\mu$, and its image passes
through the chosen very general point.  Full nef dimension and
integrality therefore give $L_0\cdot C\geq1$.  Since
$P_0-rL_0$ is nef, we obtain $P_0\cdot C\geq r$, which proves
\eqref{j:bounded-curve-obstruction}.

If $f=1$, take $C=Y$ directly.  Then
$P_0\cdot Y\geq rL_0\cdot Y\geq r$, and the required bound is
$r\leq C_0$.  Finally, \cite[Theorem~1.3]{BZ} applies to every
multiple of its prescribed integer.  A common multiple of the
finitely many integers $m_1,\ldots,m_n$ therefore gives the last
assertion.
\end{proof}

In particular, the hypotheses of Lemma~\ref{lemma:bounded-curves}
cannot hold along a sequence with $r\to\infty$ and with $C_0,C_1$
bounded, even if the varieties and divisors vary.  The following
observation explains how full nef dimension and pseudo-effectivity
of the canonical divisor pass to the moving subvarieties used below.

\begin{lemma}[Very general images]
\label{j:very-general-images}
Let $X$ be a normal integral projective variety that is not uniruled,
and let $L$ be a nef Cartier divisor of full nef dimension on $X$.
There is a complement $X^\circ$ of a countable union of proper
closed subsets of $X$ with the following property.  If $Y$ is a
smooth integral projective variety and $g:Y\to X$ is generically
finite onto a positive-dimensional image meeting $X^\circ$, then
$g^*L$ has full nef dimension, $Y$ is not uniruled, and $K_Y$ is
pseudo-effective.
\end{lemma}

\begin{proof}
Choose $X^\circ$ so that every curve through any of its points has
positive $L$-degree and no rational curve meets $X^\circ$.  The first
condition is the very general characterization of full nef
dimension.  For the second, rational curves are parametrized by
countably many algebraic families; the closure of the evaluation
image of each such family is proper because $X$ is not uniruled.

Since $g(Y)$ meets $X^\circ$, it is contained in none of the closed
sets excluded in its definition.  A very general point of $Y$ thus
maps into $X^\circ$ and lies in the quasi-finite locus of $g$.
Every curve through such a point has a curve as its image, so its
$g^*L$-degree is positive.  This proves full nef dimension.  A
covering family of rational curves on $Y$ would similarly give a
nonconstant rational curve meeting $X^\circ$, a contradiction.
The pseudo-effectivity of $K_Y$ now follows from \cite{BDPP}.
\end{proof}

\subsection{Subadjunction and high-multiplicity components}

We use the following five results from the section
\emph{Tools for moving base components} of \cite{LA}.  Their
hypotheses distinguish generic information along a center from
global information about the ambient variety.  Intersections on
an integral subvariety are computed after pullback to its
normalization and a smooth projective resolution.  A base ideal
is \emph{isolated at the generic point} of a subvariety when its
localization there is primary for the maximal ideal.

\begin{lemma}[Numerical generic subadjunction, \cite{LA}, Lemma~7.1]
\label{j:subadjunction}
Let $Z$ be a projective $\Q$-factorial klt variety, let
$\Theta\geq0$ be a rational divisor, and let $V\subset Z$ be
integral of dimension $f>0$.  Suppose that $(Z,\Theta)$ is log
canonical at the generic point of $V$ and has a divisor of log
discrepancy zero with center $V$.  If $P$ is a nef rational Cartier
divisor on $V$ and $\rho:V^*\to V$ is a smooth projective resolution
factoring through its normalization, then
\[
 K_{V^*}\cdot(\rho^*P)^{f-1}
 \leq (K_Z+\Theta)|_V\cdot P^{f-1}.
\]
\end{lemma}

Only generic log canonicity is required.  The proof in \cite{LA}
uses the dlt modification of \cite[Theorem~2.10]{FH}, whose
truncated boundary leaves an effective surplus supported over
the non-lc locus.  On a minimal stratum over $V$, adjunction
gives a pair that is klt over the generic point of its
connected-fiber base.  This suffices for the klt-trivial
canonical bundle formula and its b-nef moduli divisor
\cite[Definition~3.1 and Theorem~3.5]{FG}.  The effective
discriminant, the moduli divisor, and finite ramification have
nonnegative intersections with the pulled-back nef class.
The comparison uses canonical Weil cycles, so it requires no
$\Q$-Gorenstein assumption on the normalization of $V$.

\begin{lemma}[A base component of high multiplicity, \cite{LA}, Lemma~7.2]
\label{j:base-component}
Let $Z$ be a projective $\Q$-factorial klt variety of dimension
$d$, let $R$ be a nef rational Cartier divisor, and let $k>0$
be an integer such that $kR$ is Cartier.  Let
$0\ne\mathcal H\subset H^0(Z,\OO_Z(kR))$ be a linear subspace
with proper base locus and base ideal $\mathfrak b$.  Suppose that
$V$ is an integral base-locus component of dimension $f<d$,
isolated at its generic point $\eta_V$, and that
$\eta_V\in Z_{\mathrm{sm}}$.  Let
$\mathfrak m$ be the maximal ideal of $\OO_{Z,\eta_V}$.
If $\mathfrak b\OO_{Z,\eta_V}\subset\mathfrak m^h$ for an
integer $h>0$, there are a rational number $0<c\leq(d-f)k/h$ and an
effective rational divisor $\Theta\sim_{\Q}cR$ such that
$(Z,\Theta)$ is log canonical at $\eta_V$ and has a log canonical
place centered on $V$.  The degree and canonical bounds are
\begin{equation}
\label{j:base-bounds}
 \begin{split}
 R^f\cdot V&\leq(k/h)^{d-f}R^d,\\
 K_{V^*}\cdot R^{f-1}
   &\leq(K_Z+cR)|_V\cdot R^{f-1}\qquad(f>0),
 \end{split}
\end{equation}
where $V^*$ is a smooth projective resolution.  The degree bound
also holds for $f=0$.
\end{lemma}

The degree estimate is the local multiplicity bound $h^{d-f}$
combined with successive general cuts by the given linear series.
The boundary is obtained by averaging general members with total
coefficient equal to the generic log canonical threshold of
$\mathfrak b$.  Thus it is an actual divisor to which
Lemma~\ref{j:subadjunction} applies.

\subsection{Full jet kernels and restriction to a fiber}

\begin{lemma}[Moving multiplicity, \cite{LA}, Lemma~7.3]
\label{j:moving-kernel}
Let $Z$ be an integral projective variety of dimension $d$, let
$R$ be a nef, big, semiample rational Cartier divisor, and fix
an integer $k>0$ such that $kR$ is Cartier.  Choose positive real
numbers $\tau_0<\cdots<\tau_d$ with consecutive gap $\delta>0$.
For $x\in Z_{\mathrm{sm}}$, put
\[
 \mathcal H_i(x)=H^0\bigl(Z,\OO_Z(kR)
                  \otimes\mathfrak m_x^{\lceil k\tau_i\rceil}\bigr).
\]
Suppose that for very general $x$ all these spaces are nonzero and
the base locus of $\mathcal H_0(x)$ has a positive-dimensional
component through $x$.  If $k\delta\geq4$, then, after an
algebraic parameter extension and restriction to nonempty opens,
there are an integral parameter variety $T$ with a dominant
marked-point map $x:T\to Z_{\mathrm{sm}}$, a fixed index
$i\in\{0,\ldots,d-1\}$, and a family of integral subvarieties
$V_t$ containing $x(t)$ and common to the base loci of
$\mathcal H_i(x(t))$ and $\mathcal H_{i+1}(x(t))$.  In particular,
the incidence
\[
 \Gamma=\{(t,z):z\in V_t\}\subset T\times Z
\]
dominates $Z$.

For general $t$, the dimension $f$ of $V_t$ lies in $[1,d-1]$,
$V_t$ is an isolated component of the higher base locus at its
generic point, and the entire higher-kernel base ideal is
contained there in the ordinary power $\mathcal I_{V_t}^h$, where
\[
 h=\lceil k\delta/2\rceil,\qquad k/h\leq2/\delta.
\]
The incidence base ideal is likewise contained in
$\mathcal I_\Gamma^h$ on a dense open of $\Gamma$.
These containments are taken in the smooth ambient open.
If $Z$ is also $\Q$-factorial and klt, then
\begin{equation}
\label{j:moving-bounds}
 \begin{split}
 R^f\cdot V_t&\leq(2/\delta)^{d-f}R^d,\\
 K_{V_t^*}\cdot R^{f-1}
   &\leq(K_Z+cR)|_{V_t}\cdot R^{f-1},
       \qquad 0<c\leq2(d-f)/\delta.
 \end{split}
\end{equation}
The top intersection on $V_t$ is positive when $V_t$ meets the
open on which the big semiample contraction is an isomorphism.
\end{lemma}

The use of the full kernels matters: differentiation in the
parameter lowers vanishing at the moving mark and still gives a
global section in the lower kernel.  Incidence dominance makes
these parameter directions span the normal directions.  The
vanishing of all derivatives of order less than $h$ then gives
the ordinary-power containment, which can be restricted to
general fibers as follows.

\begin{lemma}[Restriction of an isolated base ideal, \cite{LA}, Lemma~7.4]
\label{j:fiber-restriction}
In the setting of Lemma~\ref{j:moving-kernel}, let $g:Z\to B$ be
a morphism.  There is a dense open of the incidence with the
following property.  Choose $(t,z)$ in this open, put $b=g(z)$,
and suppose that $V_t$ is smooth over a smooth open of its actual
image at $z$.  Let $F$ be the reduced component through $z$ of
the scheme fiber of $V_t\to g(V_t)$.  Near $z$ in $Z_b$, the
restricted higher-kernel base ideal has support exactly $F$ and
is contained in the ordinary power $\mathcal I_F^h$.

If $F$ is proper in an integral ambient fiber component containing
it, the restricted series is nonzero there and $F$ is an isolated
base component at its generic point.  Applying
Lemma~\ref{j:base-component} on that ambient component additionally
requires its projectivity, $\Q$-factoriality, klt singularities,
and smoothness at the generic point of $F$.
\end{lemma}

Here the fiber is chosen through a general incidence point.  The
statement does not assert the same conclusion for an arbitrary
special fiber.  Smoothness over the actual image identifies the
restricted ideal of $V_t$ with the reduced ideal of $F$ locally;
this is what preserves ordinary powers.

The $\Q$-factoriality needed for the base-component estimates can
also be supplied by a small projective $\Q$-factorialization
$\sigma:Z'\to Z$ of a klt ambient variety.  This passage is made
only for components that meet the isomorphism open of $\sigma$.
Their strict transforms retain the generic base-ideal and
multiplicity conditions for the pulled-back series, and the
testing class $\sigma^*R$ is nef.  Since
$K_{Z'}=\sigma^*K_Z$, projection formula gives the same degree
and canonical-intersection bounds on a common resolution.
The pullback testing class need not be ample; nefness is the
property used in these estimates.  Passing to higher resolutions
does not change the canonical intersections, because their
additional exceptional divisors have zero intersection with
the pulled-back testing class to the required power.

\subsection{Finite covers with a fixed branch complement}

A bound on the degree of a branched cover does not make the
possible covers finite.  The final companion input supplies the
needed fixed branch complement from a family whose branch
divisors do not sweep the base.

\begin{lemma}[Finite covers after excluding sweeping branch divisors,
              \cite{LA}, Lemma~7.5]
\label{j:finite-covers}
Let $Y$ be a normal integral projective variety, let
$\mathcal V\to T$ be a smooth projective family with integral
fibers over an integral parameter variety, and let
$g:\mathcal V\to Y$ be a morphism.  Suppose that the general
fiber maps $g_t:V_t\to Y$ are dominant and generically finite of
degree at most a fixed integer $e$.

After a finite parameter extension and shrinking, spread the
geometric generic components of the zero divisor of the relative
Jacobian over $Y_{\mathrm{sm}}$ to divisors
$\mathcal R_1,\ldots,\mathcal R_N$ in $\mathcal V$.
Retain this notation after the base change.  Let $I$ consist of
the indices $j$ for which
$\overline{g_t((\mathcal R_j)_t)}$ is a divisor in $Y$ for general
$t$.  Assume that
\[
 \overline{g(\mathcal R_j)}\ne Y\qquad(j\in I).
\]
After shrinking again, there is a smooth nonempty open
$Y^\circ\subset Y$ over which every general finite normal Stein
cover of $g_t$ is finite \'etale.  Only finitely many such normal
covers of $Y$ occur, up to isomorphism over $Y$.
The conclusion holds simultaneously for finitely many
projections; the open and the finite lists may depend on all
fixed parameters of the family.
\end{lemma}

To see the role of the hypothesis, remove $Y_{\mathrm{sing}}$
and the finitely many proper image closures
$\overline{g(\mathcal R_j)}$ for $j\in I$.  Every branch prime
of a general Stein cover is accounted for by one of these
relative Jacobian components.  Purity therefore makes the
restricted covers \'etale.  The fundamental group of the
resulting smooth complex variety is finitely generated, so it
has only finitely many covers of degree at most $e$.  Finally,
a finite normal cover of $Y$ is determined by its function
field, hence by its restriction to this fixed open.

\section{Jets on a terminal variety with trivial canonical class}
\label{sec:geometry}

The reduction to a Calabi--Yau pair leaves a positivity problem for a
nef divisor: full nef dimension must imply bigness. We first treat the
case with trivial canonical bundle and terminal singularities. The
argument follows the scalar and two-slot jet constructions of
\cite[Sections~7--8]{LA}. Its geometric exclusions differ from the
nonvanishing argument there: lower-dimensional instances of $(P_j)$
and the polarized effective birationality theorem replace the
exclusions for a canonical nonvanishing counterexample.

\begin{theorem}\label{thm:terminal-big}
Let $n\ge2$, and assume $(P_j)$ for $j<n$. Let $X$ be a projective
$\Q$-factorial terminal complex variety of dimension $n$ with
$K_X\sim0$. If $L$ is a nef Cartier divisor on $X$ with full nef
dimension, then $L$ is big.
\end{theorem}

The proof occupies this section and Section~\ref{sec:frobenius}.
Assuming that $L$ is not big, we obtain two contrasting jet estimates:
sections on $X$ have small normalized vanishing order, whereas a
projective bundle over $X\times X$ admits a large jet system. The
Frobenius comparison will show that these estimates are incompatible.

\subsection{Normalization and the curve obstruction}

Throughout this section suppose that $X,L$ form a counterexample to
Theorem~\ref{thm:terminal-big}. Fix an ample Cartier divisor $A$ on $X$
and a positive rational number $\epsilon$ such that
\begin{equation}\label{j:epsilon}
 (14\epsilon)^n<\frac14,
 \qquad 80\epsilon<\frac34.
\end{equation}
Let $t>0$ tend to zero along a sequence of rational numbers. Since
$L$ is nef and not big, $L^n=0$. We may therefore choose positive
integers $r=r(t)$ satisfying
\begin{equation}\label{eq:scaling}
 P=P_t=r(L+tA),\qquad
 r\longrightarrow\infty,\qquad P^n\longrightarrow\epsilon^n.
\end{equation}
For example, round $\epsilon/((L+tA)^n)^{1/n}$ to the nearest positive
integer. The divisor $P$ is ample and rational Cartier, and $P-rL$ is
nef. After discarding finitely many terms, we have
\begin{equation}\label{j:volume-window}
 \frac{\epsilon^n}{2}\le P^n\le2\epsilon^n.
\end{equation}
Every curve through a very general point of $X$ has positive integral
$L$-degree, hence $P$-degree at least $r$. Thus a covering family of
curves of bounded $P$-degree will give a contradiction.

We specify how the same obstruction applies to subvarieties and
covers. A resolution $W\to X$ has an effective canonical divisor:
pull back a nowhere-zero canonical form on the smooth locus of $X$
and use terminality. Hence $W$, and therefore $X$, is not uniruled
\cite{BDPP}. By Lemma~\ref{j:very-general-images}, if a smooth
projective variety $Y$ maps generically finitely to a
positive-dimensional subvariety containing a sufficiently very general
point of $X$, then $Y$ is not uniruled and the pulled-back Cartier
divisor $L_Y$ has full nef dimension. In particular, $K_Y$ is
pseudo-effective. If $f=\dim Y<n$, the induction hypothesis $(P_f)$
applied to $(Y,0)$ gives
\begin{equation}\label{j:lower-big}
                         K_Y+L_Y\ \text{big}.
\end{equation}
All occurrences of ``very general'' below include the fixed
conditions of Lemma~\ref{j:very-general-images}.

Consequently, for $f<n$ it suffices to find on such a $Y$ a nef
testing divisor $P_Y$ with
\[
 0<P_Y^f\le C_0,\qquad P_Y-rL_Y\ \text{nef},\qquad
 K_YP_Y^{f-1}\le C P_Y^f,
\]
where $C_0,C$ are independent of $t$. Indeed
$L_YP_Y^{f-1}\le P_Y^f/r$, so Lemma~\ref{lemma:bounded-curves}
applies. We will also use that lemma in dimension $n$, after proving
the required bigness separately. For the rest of this section, a
bound is called uniform if it is independent of $t$; it may depend
on parameters fixed before $t$ tends to zero. Since the dimensions
lie in a finite range, Lemma~\ref{lemma:bounded-curves} excludes all
such configurations once $r$ is sufficiently large.

\subsection{The scalar jet estimate}

Put $\rho=(P^n)^{1/n}$. The first estimate bounds the vanishing order
of every section at a very general point; it is stronger than a
bound on a single fixed linear system.

\begin{lemma}\label{lemma:scalar-bound}
In the setting of \eqref{eq:scaling}, for all sufficiently small $t$
and at a very general point $x\in X$, every nonzero section of $kP$,
where $k$ is a positive integer and $kP$ is Cartier, satisfies
\begin{equation}\label{j:scalar-order}
                         \ord_x(s)\le4k\rho.
\end{equation}
\end{lemma}

\begin{proof}
Suppose that violations occur for arbitrarily small $t$. Fix such a
$t$. At a very general point a violation for some $k$ implies that
the corresponding complete jet kernel is nonzero at the general
point. Indeed the kernel ranks are algebraic rank conditions on the
smooth locus, and there are only countably many choices of $k$ and
of the required integral order. Raising the violating section to
powers allows $k$ to be arbitrarily large and divisible.

Choose $n+1$ levels
\[
 \tau_i=\frac32\rho+i\delta\quad(0\le i\le n),
 \qquad \delta=\frac{2\rho}{n}.
\]
They all lie strictly between $\rho$ and $4\rho$. For a sufficiently
large divisible $k$, their full jet kernels are nonzero and
$k\delta\ge4$. The lowest kernel has a positive-dimensional base
component through its marked point. Otherwise that point would be
an isolated base component of multiplicity at least
$h_0=\lceil k\tau_0\rceil$. The zero-dimensional case of
Lemma~\ref{j:base-component} would give
\[
                1\le (k/h_0)^nP^n<1,
\]
a contradiction.

Lemma~\ref{j:moving-kernel} now supplies a moving family of proper
base components $V$, with dominant incidence and dimension
$1\le f<n$. On a resolution $V^*$, write $P_V,L_V$ for the
restrictions followed by pullback. The lemma gives
\[
 0<P_V^f\le(2/\delta)^{n-f}P^n,
 \qquad K_{V^*}P_V^{f-1}\le cP_V^f,
 \qquad 0<c\le\frac{2(n-f)}{\delta}.
\]
Here the canonical term of the ambient variety vanishes because
$K_X\sim0$. The volume window \eqref{j:volume-window} bounds
$\delta$ above and away from zero, so these are uniform estimates.
The incidence point may be chosen over a very general point of $X$.
Thus \eqref{j:lower-big} gives bigness of $K_{V^*}+L_V$, and
$P_V-rL_V$ is nef. Lemma~\ref{lemma:bounded-curves} contradicts
$r\to\infty$.
\end{proof}

\subsection{The projective bundle and its volume}

We now seek a jet system whose order is much larger than the scalar
bound. The additional direction comes from the fibres of a
projective bundle; its two summands keep track of the two copies of
$X$. For a divisor on $X$, subscripts $1,2$ denote pullback from the
two factors of $X\times X$. Set
\begin{equation}\label{eq:bundle}
 Z=\PP\bigl(\OO_X(L)_1\oplus\OO_X(L)_2\bigr)
       \xrightarrow{\pi}X\times X,
 \qquad \xi=c_1(\OO_Z(1)),
 \qquad R=P_1+P_2+q\xi,
\end{equation}
where $q$ is a positive integer to be fixed shortly. We use the
convention
$\pi_*\OO_Z(a)=\operatorname{Sym}^a(\OO_X(L)_1\oplus\OO_X(L)_2)$
for $a\ge0$. The tautological class $\xi$ is nef, and $R$ is ample.
The product and projective-bundle descriptions show that $Z$ is klt.
The two natural sections, called the axes, have classes
$\xi-L_1$ and $\xi-L_2$. The canonical bundle formula gives
\begin{equation}\label{j:axes-canonical}
                         K_Z\sim-\text{(sum of the two axes)}.
\end{equation}
We call the complement of the axes the torus open.

The fibre of either projection $Z\to X$ over a smooth point is,
up to a constant one-dimensional factor,
\[
 Z_{\mathrm{sl}}=\PP(\OO_X\oplus\OO_X(L)),
 \qquad R_{\mathrm{sl}}=P+q\xi.
\]
Its axes have classes $\xi$ and $\xi-L$, and its canonical divisor
is again minus their sum. The slice is projective and klt. In
particular it has the singularity hypotheses needed to apply the
base-component estimates, using a small $\Q$-factorialization if
necessary as explained after Lemma~\ref{j:moving-kernel}.

Write $d=2n+1$. There are positive constants depending only on $n$
such that, for each fixed $q$ and all sufficiently small $t$,
\begin{equation}\label{j:bundle-volumes}
 c_nq\epsilon^{2n}\le R^d\le C_nq\epsilon^{2n},
 \qquad R_{\mathrm{sl}}^{n+1}\le C_{\mathrm{sl},n}q\epsilon^n.
\end{equation}
To see this, expand in nef classes and use
\[
 \pi_*(\xi^j)=\sum_{a=0}^{j-1}L_1^aL_2^{j-1-a}
 \quad(j\ge1).
\]
The term containing exactly one $\xi$ is
$(2n+1)\binom{2n}{n}q(P^n)^2$, giving the lower bound.
For the upper bound, every occurrence of $L_i$ can be bounded in
mixed nef intersections by $P_i/r$. After factoring out $q$, the
remaining powers of $q/r$ are bounded by one once $r\ge q$.
There are only finitely many terms depending on $n$. The slice
calculation is the same, using $\pi_*(\xi^j)=L^{j-1}$.
Together with \eqref{j:volume-window}, this proves
\eqref{j:bundle-volumes}.

For a nef rational Cartier divisor $R$ and a smooth point $z$, its
\emph{Seshadri constant} is
\[
 \varepsilon(R;z)=\inf_{C\ni z}\frac{R\cdot C}{\mult_z C},
\]
where $C$ ranges over integral curves through $z$. Thus a lower bound
for this constant controls degree relative to multiplicity at the point.

\begin{proposition}\label{prop:large-seshadri}
For a suitable fixed positive integer $q$, and all sufficiently
small $t$ in \eqref{eq:scaling}, there is a smooth point $z\in Z$
in the torus open, lying over smooth points of both factors, such
that
\begin{equation}\label{j:large-seshadri}
                         \varepsilon(R;z)>2n+2.
\end{equation}
More generally, for any prescribed dense open subsets $U_1,U_2$
of the smooth locus of $X$, the point may be chosen over
$U_1\times U_2$. In particular, these opens may be the
isomorphism loci of fixed resolutions of $X$.
\end{proposition}

We prove the proposition in the remainder of this section. Choose a
positive rational gap $\delta$ such that
\begin{equation}\label{j:slice-gap}
             (2/\delta)^nC_{\mathrm{sl},n}\epsilon^n<1.
\end{equation}
Next choose $q$ so large that $d+1$ increasing levels of gap $\delta$,
with lowest level strictly greater than $2n+2$, lie strictly below
$(c_nq\epsilon^{2n})^{1/d}$. These choices precede the limit
$t\to0$.

Fix the prescribed opens $U_1,U_2$, if any. Suppose that
\eqref{j:large-seshadri} fails throughout their smooth torus
preimage for arbitrarily small $t$. For any such $t$, asymptotic
Riemann--Roch and the elementary
jet count at a smooth point show that all the full kernels at the
chosen levels are nonzero for sufficiently large divisible $k$.
Indeed the highest level has $d$-th power smaller than $R^d$.
Failure of \eqref{j:large-seshadri} supplies, through each very
general point, a curve with degree divided by multiplicity smaller
than the lowest level. Every section in the lowest kernel vanishes
identically on that curve: otherwise its intersection with the
curve is at least its vanishing order times the curve's
multiplicity, contradicting the degree inequality. Thus the lowest
base locus has a positive-dimensional component through the mark.
Take $k\delta\ge4$ and apply Lemma~\ref{j:moving-kernel}.

We obtain a family of proper integral components $V\subset Z$ whose
incidence dominates $Z$. For its general members, of dimension $f$,
the estimates are
\begin{equation}\label{j:total-component-bounds}
 \begin{split}
 0<R^f\cdot V&\le(2/\delta)^{d-f}R^d,\\
 K_{V^*}R^{f-1}&\le(K_Z+cR)|_V R^{f-1},
 \qquad 0<c\le2(d-f)/\delta.
 \end{split}
\end{equation}
Here and below classes on a resolution are understood by pullback.
The positivity follows from ampleness of $R$. General members meet
the torus open, so \eqref{j:axes-canonical} makes the ambient
canonical contribution nonpositive. All constants in these bounds
are uniform now that $q,\epsilon,\delta$ have been fixed.

\subsection{Excluding positive-dimensional slice fibres}

We first constrain the two projections of $V$ to $X$. For either
projection, consider the component $F$ of a general fibre of
$V\to\operatorname{im}(V\to X)$ through a general incidence point.
Choose that incidence point in the opens of
Lemma~\ref{j:fiber-restriction}, over very general points of both
factors of $X$, and in the smooth torus open. We claim that
\begin{equation}\label{j:no-intermediate-fibre}
                         \dim F\notin\{1,\ldots,n\}.
\end{equation}

Suppose instead that $1\le f_{\mathrm{sl}}=\dim F\le n$.
The ambient fibre is $Z_{\mathrm{sl}}$, of dimension $n+1$.
The restriction lemma gives a nonzero restricted subseries with
$F$ an isolated base component and base ideal contained generically
in $\mathcal I_F^h$, where $h=\lceil k\delta/2\rceil$.
Generic smoothness is imposed over the actual image of $V$; this is
why a general fibre through the incidence point was chosen.
The point is also smooth in the ambient slice. Applying
Lemma~\ref{j:base-component} on that slice gives
\begin{equation}\label{j:slice-bounds}
 \begin{split}
 0<R_{\mathrm{sl}}^{f_{\mathrm{sl}}}\cdot F
   &\le(2/\delta)^{n+1-f_{\mathrm{sl}}}
                      R_{\mathrm{sl}}^{n+1},\\
 K_{F^*}R_{\mathrm{sl}}^{f_{\mathrm{sl}}-1}
   &\le(K_{Z_{\mathrm{sl}}}+cR_{\mathrm{sl}})|_F
                      R_{\mathrm{sl}}^{f_{\mathrm{sl}}-1},
 \end{split}
\end{equation}
with $c\le2(n+1-f_{\mathrm{sl}})/\delta$. Cartier multiples are
chosen before restriction, so the constant line from the fixed
factor causes no change in this calculation.

Let $W_0\subset X$ be the image of $F$ in the other factor.
There are three cases: $F$ is generically finite over a proper
subvariety of $X$; $F$ contains the projective-line fibres over its
image; or $F$ is a multisection over all of $X$.

\paragraph{A generically finite map to a proper subvariety.}
If $F\to W_0$ is generically finite and $\dim W_0<n$, use
$P_Y=R_{\mathrm{sl}}|_{F^*}$ and the pullback $L_Y$ from $X$.
The first inequality in \eqref{j:slice-bounds} gives a uniform
positive volume bound. Since $F$ meets neither axis generically,
its ambient canonical term is nonpositive. Hence the second
inequality gives a uniform canonical-intersection ratio.
Moreover $P_Y-rL_Y$ is nef. The image $W_0$ contains the very
general projection of the incidence point, so
\eqref{j:lower-big} applies. Lemma~\ref{lemma:bounded-curves}
excludes this case for large $r$.

\paragraph{A projective-line bundle over a proper subvariety.}
If $F\to W_0$ is not generically finite, its general fibre has
dimension one. Since the ambient projection is a projective-line
bundle and $F$ is integral, $F$ is the full bundle over $W_0$.
Write $w=\dim W_0=f_{\mathrm{sl}}-1$.
If $w=0$, then $F$ is a projective-line fibre and
$R_{\mathrm{sl}}\cdot F=q$. The first inequality in
\eqref{j:slice-bounds}, together with
\eqref{j:bundle-volumes} and \eqref{j:slice-gap}, gives
\[
 q\le(2/\delta)^nR_{\mathrm{sl}}^{n+1}
   \le(2/\delta)^n C_{\mathrm{sl},n}q\epsilon^n<q,
\]
which is impossible.

Suppose $w>0$. Resolve $W_0$ to $W_0^*$ and take as a resolution
of $F$ the induced projective bundle
$p_0:F^*\to W_0^*$. Its relative canonical divisor is
$-2\xi+p_0^*L$, exactly the restriction of
$K_{Z_{\mathrm{sl}}}$. Cancelling this term in
\eqref{j:slice-bounds} gives
\begin{equation}\label{j:saturated-canonical}
 p_0^*K_{W_0^*}R_{\mathrm{sl}}^w
                   \le cR_{\mathrm{sl}}^{w+1}\cdot F.
\end{equation}
The divisor $K_{W_0^*}$ is pseudo-effective by the very general
point condition. On expanding the left side and pushing down, its
term with exactly one $\xi$ is
$wqK_{W_0^*}P^{w-1}$; every other nonzero term is nonnegative,
since it pairs the pseudo-effective canonical divisor with nef
classes. On the right, the same expansion used for
\eqref{j:bundle-volumes} gives
\begin{equation}\label{j:saturated-volume}
 (w+1)qP^w\cdot W_0
 \le R_{\mathrm{sl}}^{w+1}\cdot F
 \le C'_nqP^w\cdot W_0,
\end{equation}
using $r\ge q$. The slice degree bound and the first inequality in
\eqref{j:saturated-volume} bound $0<P^w\cdot W_0$ uniformly.
Equations \eqref{j:saturated-canonical}--\eqref{j:saturated-volume}
also give
\[
 K_{W_0^*}P^{w-1}\le\frac{cC'_n}{w}P^w\cdot W_0.
\]
Apply \eqref{j:lower-big} and Lemma~\ref{lemma:bounded-curves}
on $W_0^*$ with the pullbacks of $P,L$. This excludes the second
case.

\paragraph{A multisection over $X$.}
The only remaining possibility is that $f_{\mathrm{sl}}=n$ and
$F\to X$ is generically finite of degree $e>0$. Intersect $F$ with
each of the two Cartier axes and push the resulting cycles down
to $X$. This gives effective integral Weil divisors $D_1,D_2$.
The axes do not contain $F$, and their classes differ by $L$;
projection of the associated rational section, or equivalently its
norm, gives
\begin{equation}\label{j:signed-L}
                         D_1-D_2\sim\pm eL.
\end{equation}
For either axis $S$, the classes $R_{\mathrm{sl}}-qS$ and
$R_{\mathrm{sl}}-P$ are nef. Intersecting with the effective cycle
$S\cdot F$ and then using nefness on $F$ yields
\begin{equation}\label{j:axis-degrees}
 qD_iP^{n-1}\le qS\cdot F\cdot R_{\mathrm{sl}}^{n-1}
                  \le R_{\mathrm{sl}}^n\cdot F.
\end{equation}
Thus both axis-image degrees are uniformly bounded.

Set $G=D_1+D_2$. It is $\Q$-Cartier because $X$ is
$\Q$-factorial. Choose $b>0$ rational sufficiently small that
$(X,bG)$ is klt. We claim
\begin{equation}\label{j:positive-boundary-kappa}
                         \kappa(K_X+bG)>0.
\end{equation}
The Iitaka dimension is at least zero, since
$K_X+bG\sim_\Q bG$. If it were zero, Theorems~\ref{input:models}
and \ref{input:abundance} would give a finite ordinary MMP to a
good model $Y$ with transformed adjoint rationally linearly
trivial. The transform of $K_X$ is linearly trivial as well.
Consequently the effective transform of $G$ is zero: all its
components have been contracted. On a common resolution
$a:U\to X$, $b_Y:U\to Y$, relation \eqref{j:signed-L} then
expresses $a^*L$ up to rational linear equivalence as a signed
$b_Y$-exceptional divisor $E$. Divisors exceptional over $X$ are
also exceptional over $Y$, since the MMP extracts none.
Nefness of $a^*L$ and the negativity lemma imply $E\le0$.
Its class is also pseudo-effective. A nonzero anti-effective
divisor has negative intersection with a sufficiently high power
of an ample divisor, so $E=0$. This contradicts full nef dimension
of $L$, and proves \eqref{j:positive-boundary-kappa}.

The already proved positive-Iitaka-dimension case
(Proposition~\ref{prop:full-positive}) now gives
\[
                         K_X+bG+L\ \text{big}.
\]
Its full-positivity hypothesis holds because $L$ has full nef
dimension and $K_X+bG$ is pseudo-effective.
Take a log resolution $a:Y\to X$ of $G$ and let $D_0$ be the
reduced simple normal crossings divisor consisting of the strict
support of $G$ and every exceptional divisor. Decrease $b$ if
necessary; Proposition~\ref{prop:full-positive} still applies.
Terminality and the support of $a^*G$ then give
\[
 K_Y+D_0+a^*L\ \ge\ a^*(K_X+bG+L)
\]
up to rational linear equivalence and an effective difference.
Indeed the finitely many exceptional discrepancies are positive,
and decreasing $b$ makes each coefficient of $b\,a^*G$ no larger
than the corresponding coefficient of $K_Y-a^*K_X+D_0$.
Hence the left side is big.

Test on $Y$ with $a^*P$. Exceptional divisors contribute zero to
intersection with $(a^*P)^{n-1}$, and $K_X\sim0$. The degree of
the reduced strict support is no larger than $GP^{n-1}$, bounded
by \eqref{j:axis-degrees}; also
$LP^{n-1}\le P^n/r$. Since \eqref{j:volume-window} bounds $P^n$
away from zero, we obtain a uniform constant $C$ such that
\[
 (K_Y+D_0+a^*L)(a^*P)^{n-1}\le C(a^*P)^n.
\]
Lemma~\ref{lemma:bounded-curves}, now with its reduced boundary
$D_0$, gives the final contradiction. This proves
\eqref{j:no-intermediate-fibre}.

Thus a general moving component $V$ has fibre dimension either
zero or $n+1$ over its image in either factor. Fibre dimension
$n+1$ would make $V$ the full inverse image of its image in that
factor. Since $V$ is proper, the image has dimension at most
$n-1$, so the other projection would have fibre dimension between
$1$ and $n$, already excluded. Both projections of $V$ are
therefore generically finite onto their images, and $\dim V\le n$.
If $\dim V<n$, the bounds \eqref{j:total-component-bounds}, the
nef class $R-rL_i$, and \eqref{j:lower-big} again contradict
Lemma~\ref{lemma:bounded-curves}. Only $n$-dimensional
correspondences, dominant over both factors, remain.

\subsection{Correspondences and their branch divisors}

For these remaining $V$, \eqref{j:total-component-bounds} gives
uniform bounds
\begin{equation}\label{j:correspondence-bounds}
 R^n\cdot V\le C,
 \qquad K_{V^*}R^{n-1}\le cR^n\cdot V\le C,
\end{equation}
after increasing $C$. Each projection has bounded degree:
$R-P_i$ is nef, and $P^n$ is bounded below.

For each fixed $t$, resolve the general members in their algebraic
family and shrink the parameter space so that the resolutions form
a smooth projective family with integral fibres and morphisms to
$X$ in both slots. The dominant incidence and all numerical bounds
persist. Finite extensions of the parameter space are allowed.
Over $X_{\mathrm{sm}}$, name the components of the relative
Jacobian divisors as in Lemma~\ref{j:finite-covers}. Consider a
component whose image on a general member is a divisor $J\subset X$.

Such divisors $J$ cannot sweep $X$. If they did, choose a general
one through a very general point. Let $E_J$ be a ramification
divisor over it on the corresponding $V^*$. Pullback of a
canonical form from $X$ gives an effective canonical divisor on
$V^*$. Regularity follows by factoring through a resolution of
$X$ and using its canonical singularities; the coefficient of
$E_J$ is positive because its image meets $X_{\mathrm{sm}}$
generically. The difference $R-P_i$ is nef, so
\eqref{j:correspondence-bounds} and projection formula bound
\[
                    0<d_J:=P^{n-1}\cdot J\le C.
\]
On a resolution $J^*$, Lemma~\ref{j:subadjunction}, applied to
$(X,J)$ at the generic point of $J$, gives
\begin{equation}\label{j:branch-canonical}
 \begin{split}
 K_{J^*}P^{n-2}
 &\le(K_X+J)\cdot J\cdot P^{n-2}\\
 &=J^2P^{n-2}
 \le\frac{(JP^{n-1})^2}{P^n}
 \le C' d_J.
 \end{split}
\end{equation}
The pair is lc at that generic point because $X$ is normal and
hence smooth there. The middle inequality is the Hodge index
inequality, applied on a resolution and obtained by ample
approximation of the pulled-back testing class. It does not require
$J$ to be nef. The final constant is uniform by the volume lower
bound and the bound for $d_J$.

Since $J$ sweeps $X$, its resolution satisfies the very general
point conditions used in \eqref{j:lower-big}. With testing class
$P|_{J^*}$, whose top self-intersection is $d_J$, Equation
\eqref{j:branch-canonical} gives precisely the hypotheses of
Lemma~\ref{lemma:bounded-curves} in dimension $n-1$. This
contradiction excludes sweeping branch divisors.

Lemma~\ref{j:finite-covers} now gives, for each slot and this fixed
$t$, only finitely many possible finite normal Stein covers of
$X$, up to isomorphism over $X$. The point of the preceding step
is that degree bounds alone would not give this conclusion:
non-sweeping branch images provide a fixed smooth branch complement
for the family. Neither that complement nor the resulting finite
lists need be uniform in $t$.

\subsection{Fixed covers and an abelian model}

Each $V^*$ maps birationally to both of its normal Stein covers,
and hence determines a graph of a birational map between a pair
of covers in the finite lists. We explain why some such fixed pair
has a family of graphs dominating its product. Restrict the finite
set of pairs to those realized by at least one original member
$V^*$ satisfying \eqref{j:correspondence-bounds}. The original
incidence dominates $Z$, so the projected incidences for these
pairs cover a dense open of $X\times X$. For a fixed pair, graphs of
birational maps lie in countably many Hilbert schemes. Stratify
these schemes so that the fibres are integral and flat and both
projections are birational; the latter conditions may be tested
by dominance and degree one after shrinking. Because $\C$ is
uncountable, a dense open of $X\times X$ cannot be covered by
countably many proper closed subsets. Thus, for one fixed pair
$T_1,T_2$, some integral graph family has dominant evaluation to
$X\times X$. Since $T_1\times T_2$ is finite over $X\times X$
and has the same dimension, this evaluation is dominant onto
$T_1\times T_2$ as well.

Both covers are non-uniruled because they are finite over $X$.
The birational-group theorem of Hanamura
\cite[Theorems~2.1--2.2]{Hanamura}, in the form recalled in
\cite[Proposition~3.7 and Theorems~3.8--3.9]{Blanc}, implies that
they are birational to an abelian variety. Here are the details
of the application. Identify the two covers birationally by one
map in the family. On a suitable smooth projective birational
model of the resulting non-uniruled variety, the reduced
birational group is a group scheme locally of finite type, with
identity component an abelian variety acting regularly. Conjugate
the graph family to this model and shrink again to obtain a flat
family of integral graphs. By the flat-graph representability
statement, its reduced integral parameter space maps to the
reduced birational group. Its connected image lies in one
translate of the identity component. Product-dominant evaluation
therefore gives a dense orbit for that abelian identity component.
The orbit is a quotient of an abelian variety by a stabilizer,
hence is itself an abelian variety. The fixed cover is birational
to it. This is the birational-group part of the fixed-cover
argument in \cite[Proposition~6.11]{LA}; no canonical
nonvanishing conclusion from that proposition is being used.

This structural conclusion depends only on the two covers. We may
therefore return to an original member $V^*$ belonging to this pair,
which retains \eqref{j:correspondence-bounds}, and resolve
further to obtain a birational morphism
\[
                         p_0:V^*\longrightarrow A_0
\]
to an abelian variety. Higher resolutions do not change the
canonical intersection with the pulled-back $(n-1)$-st power of
$R$. Let $L_0$ be the pullback of $L$ from either slot. It is nef,
Cartier, and of full nef dimension. We claim
\begin{equation}\label{j:abelian-big}
                         K_{V^*}+L_0\ \text{big}.
\end{equation}

The pushforward $N_0=(p_0)_*L_0$ is pseudo-effective on $A_0$,
and is therefore nef. Indeed translations make the
pseudo-effective and nef cones of an abelian variety coincide.
If $N_0$ were not big, the standard description of nef line
bundles on abelian varieties would give a positive-dimensional
abelian subvariety on whose translates $N_0$ is numerically
trivial. One can see this from the semipositive Hermitian form of
the numerical class: its kernel is rational for the lattice,
since the alternating form of the integral Cartier class is
integral, and therefore defines an abelian subvariety.
A general translate meets the image of the exceptional locus of
$p_0$ in codimension at least two, if at all. Complete-intersection
curves in that translate can thus be chosen through very general
points and avoiding this image. Their lifts have $L_0$-degree
zero, contradicting full nef dimension. Hence $N_0$ is big.

By the negativity lemma,
\[
                         p_0^*N_0-L_0=E_0\ge0
\]
is exceptional. The canonical divisor of $V^*$ has an effective
representative $E_K$ with positive coefficient along every
$p_0$-exceptional divisor, because the target is smooth with
trivial canonical bundle. Choose a rational $0<u<1$ sufficiently
small that $E_K-uE_0\ge0$. Then
\[
 K_{V^*}+uL_0\sim_\Q u p_0^*N_0+(E_K-uE_0)
\]
is big. Adding the nef divisor $(1-u)L_0$ proves
\eqref{j:abelian-big}.

We may now apply Lemma~\ref{lemma:bounded-curves} in dimension
$n$ on $V^*$, with testing class $R$ and empty boundary.
The bounds \eqref{j:correspondence-bounds} give both the top
intersection bound and the canonical ratio, while $R-rL_0$ is
nef and \eqref{j:abelian-big} supplies bigness. This final
contradiction excludes the correspondence case.

\paragraph{Conclusion of the proof of
Proposition~\ref{prop:large-seshadri}.}
The alternatives for every moving base component have been
exhausted: first its slice fibres, then components of dimension
less than $n$, sweeping branch divisors, and finally the fixed
covers. In each case the contradiction comes from
Lemma~\ref{lemma:bounded-curves} with constants uniform in $t$.
The finite cover lists were used only for one fixed $t$, and their
sizes do not enter the degree bounds. Thus failures of
\eqref{j:large-seshadri} cannot occur along a sequence $t\to0$.
This proves the proposition, including the additional open
conditions on the chosen point.\hfill$\square$

We have now constructed, under the hypothetical failure of
Theorem~\ref{thm:terminal-big}, the scalar estimate of
Lemma~\ref{lemma:scalar-bound} and the projective-bundle estimate
of Proposition~\ref{prop:large-seshadri}. The next section fixes
one sufficiently small value of $t$ and places these data in the
finite-data Frobenius comparison.

\section{Frobenius comparison}
\label{sec:frobenius}

The scalar estimate of Lemma~\ref{lemma:scalar-bound} and the
projective-bundle estimate of Proposition~\ref{prop:large-seshadri}
will now give the contradiction needed for
Theorem~\ref{thm:terminal-big}.  Their incompatibility is expressed
by the following theorem about fixed complex data.  It is the
finite-data Frobenius theorem of \cite[Theorem~9.1]{LA}.
We give its proof, including the bounds that
must be uniform when the residual characteristic varies.

The argument compares two orders of vanishing of one determinant.
Ordinary jets make the determinant nonzero after reduction to positive
characteristic. A small polarization bounds the rank on the diagonal,
forcing high vanishing at each diagonal point. A fixed movable curve on a point
blowup bounds the vanishing of every section of each determinant
factor and gives the opposite inequality.

\subsection{The finite-data statement}

For a vector bundle on a smooth projective curve, its slope is
degree divided by rank. The least and greatest Harder--Narasimhan
slopes are denoted by $\mu_{\min}$ and $\mu_{\max}$.
For a line bundle on a smooth variety, to generate jets through
order $a$ at a point means surjectivity to its stalk modulo
the $(a+1)$-st power of the maximal ideal.

\begin{theorem}[Finite-data Frobenius incompatibility]
\label{thm:finite-data}
Let $W$ be a smooth connected projective complex variety of
dimension $n\ge2$. Let $D,H$ be rational Cartier divisors, with
$H$ ample, let $S$ be an integral Cartier divisor, and fix an
integer $q>0$ and real numbers $r>1$, $\epsilon>0$. Suppose
\begin{equation}\label{fr:intersection-bounds}
 \begin{gathered}
 H^n\le(2\epsilon)^n,\qquad
 K_WH^{n-1}\le2H^n/r,\qquad
 (2D+qS)H^{n-1}\le4H^n,\\
 (14\epsilon)^n<\tfrac14,\qquad 80\epsilon<\tfrac34.
 \end{gathered}
\end{equation}
Assume that the following fixed data exist.
\begin{enumerate}
\renewcommand{\labelenumi}{\textup{(\roman{enumi})}}
\item A smooth integral complete-intersection flag from $W$ to
a curve $C$, whose successive divisor classes are $h_i=l_iH$
for fixed positive integers $l_i$ for which $l_iH$ is Cartier, with
$\mu_{\min}(\Omega_W^1|_C)\ge0$.
\item On the projective bundle
\[
 Z=\PP\bigl(\mathcal O_W(S)_1\oplus
                         \mathcal O_W(S)_2\bigr)\longrightarrow W\times W,
 \qquad \xi=c_1(\mathcal O_Z(1)),\qquad M=D_1+D_2+q\xi,
\]
use the convention that the direct image of $\mathcal O_Z(a\xi)$
is the $a$-th symmetric power of the displayed sum for $a\ge0$;
subscripts indicate pullback from the corresponding factor of
$W\times W$. The two axes are the sections defined by the two
summands. A smooth point $z_*$ in their complement and finitely many
sections of $k_0M$, for some integer $k_0>0$ with $k_0D$ Cartier,
generate jets through order $(2n+2)k_0$ at $z_*$.
\item A point $x\in W$, its blowup
$\pi_x:\widehat W\to W$ with exceptional divisor $J$, and a
curve class
\[
 \gamma=f_*(A_1\cdots A_{n-1}),
\]
where $f:V\to\widehat W$ is one fixed birational morphism,
$V$ is smooth integral projective, and the $A_i$ are ample
integral Cartier divisors, such that $J\cdot\gamma>0$ and
\begin{equation}\label{fr:movable-input}
 \pi_x^*(D+K_W/2+\lambda S)\cdot\gamma
 \le40\epsilon J\cdot\gamma\qquad(0\le\lambda\le q).
\end{equation}
It suffices to check the two endpoints in this affine inequality.
\end{enumerate}
These data cannot coexist.
\end{theorem}

The point $x$ is independent of the two base coordinates of $z_*$.
No nefness is required of $D,S$, or $K_W$, and the theorem does not
assume that $K_W$ is pseudo-effective. Its hypotheses specify the
flag and the movable curve class themselves. All these data will
be fixed before the residual characteristic tends to infinity.

\subsection{Reduction and Frobenius jets}
\label{fr:reduction}

We begin the proof of Theorem~\ref{thm:finite-data} by spreading
its hypotheses and converting the fixed ordinary jets
into jets of order proportional to the residual characteristic.
Put
$\Lambda=\prod_{i=1}^{n-1}l_i$. The flag satisfies
\begin{equation}\label{fr:flag-degrees}
 [C]=\Lambda H^{n-1},\qquad h_i\cdot C=l_i\Lambda H^n.
\end{equation}
Choose one sufficiently ample integral Cartier divisor $T_0$ so
that every $T_0+aS$, for integers $0\le a\le(k_0-1)q$, has a section
nonvanishing at each of the two base coordinates of $z_*$.
Fix these finitely many sections. All models, morphisms, divisors,
points, the flag, the complete-intersection witness for $\gamma$,
and the jet and filler sections spread over an integral finitely
generated $\Z$-algebra of characteristic zero. Shrink its
spectrum so that the fibers and flag are smooth projective and
geometrically integral, the birational morphisms remain birational,
the fixed ample bundles remain ample, and the finite jet
surjection and nonvanishings persist. Birationality is preserved
by spreading an isomorphism between dense opens. All displayed
intersection numbers are constant.

We also preserve $\mu_{\min}(\Omega_W^1|_C)\ge0$.
One way is to spread its characteristic-zero Harder--Narasimhan
filtration, make all its factors locally free, and use openness
of their semistability on curves; their degrees and ranks are
fixed. Equivalently, a single relative ample twist globally
generates the curve bundle. Every quotient of rank $a$ then has
degree bounded below by $-ad$ for a fixed integer $d$.
A negative-degree quotient therefore has one of finitely many ranks
and degrees. The proper relative Quot schemes for those Hilbert
polynomials have images missing the generic point: after removal
of torsion, a negative-degree coherent quotient there would
give a negative-degree vector-bundle quotient. Removing those
finitely many images proves the same openness assertion directly.

This base has closed points in arbitrarily large characteristics.
Take algebraic closures of their residue fields and exclude
denominator primes. Retain the notation for the reductions.
The curve $\gamma$ continues to pair nonnegatively with every
effective divisor: pull that divisor back under the fixed
birational morphism and intersect with its fixed ample divisors.
This tests effective divisors on the reduction itself. No
specialization of a characteristic-zero effective cone is used.

\begin{lemma}\label{fr:positive-characteristic-jets}
For a sufficiently large residual characteristic $p$, set
\[
 N_p=\lfloor p/k_0\rfloor,\qquad B_p=k_0N_pD+T_0.
\]
The sections of $pq\xi+B_{p,1}+B_{p,2}$ generate jets through order
$(2n+2)k_0N_p$ at $z_*$.  In particular, they surject onto the quotient
of its local ring by the $p$-th powers of all regular parameters.
\end{lemma}

\begin{proof}
Multiply $N_p$ copies of the fixed jet system.  Every monomial of
total degree at most $(2n+2)k_0N_p$ is a product of $N_p$ monomials
of degree at most $(2n+2)k_0$.  Taking sections with those leading
monomials gives a triangular system, ordered by total degree.
Starting at degree zero and removing the error in each successive
degree proves the required jet surjection. This argument takes
place in the local ring and never divides by factorials.

Put $d_p=(p-k_0N_p)q$, so $0\le d_p\le(k_0-1)q$.
Among the sections fixed before reduction, choose
\[
 u_0\in H^0(W,\OO_W(T_0)),\qquad
 v_{d_p}\in H^0(W,\OO_W(T_0+d_pS))
\]
nonzero at the first and second base coordinates of $z_*$,
respectively. Their tensor product belongs to the summand of
first-slot weight zero in the projective-bundle formula; its
fiber monomial is the $d_p$-th power of the second coordinate.
It defines a section of
$d_p\xi+T_{0,1}+T_{0,2}$ nonzero at $z_*$, since $z_*$ lies
off both axes. Multiplying the product jet system by this section
changes its divisor from $N_pk_0M$ to
$pq\xi+B_{p,1}+B_{p,2}$. Multiplication is invertible on the
local jet algebra. Only the finitely many previously fixed
values of $d_p$ occur.

Since $\dim Z=2n+1$, the
quotient by the $p$-th powers of regular parameters has top total
degree $(2n+1)(p-1)$.  For large $p$ this is at most
$(2n+2)k_0N_p$. The underlying ordinary/Frobenius ideal comparison is
also recorded in \cite[proof of Proposition~2.12, Equation~(2.8)]{MustataSchwede}.
\end{proof}

\subsection{Full rank from the two-slot jets}
\label{fr:full-rank}

Let $F:W\to W'$ be relative Frobenius to the base-field twist.
It is finite flat because $W$ is smooth.  Write $S'$ for the twist
of $S$, so $F^*S'=pS$, and define
\[
 \mathcal E=F_*\OO_W(B_p),\qquad s=\rk\mathcal E=p^n,
 \qquad R=s^2.
\]
For integers $0\le a\le q$, put
$U_a=H^0(W,\OO_W(B_p+paS))$.
Projection formula gives an evaluation map on $W'\times W'$:
\begin{equation}\label{fr:two-slot-evaluation}
 \bigoplus_{\substack{a+b=q\\a,b\ge0}}
 U_a\otimes U_b\otimes
 \OO(-aS'_1-bS'_2)
 \longrightarrow \mathcal E\boxtimes\mathcal E.
\end{equation}

\begin{lemma}\label{fr:generic-rank}
The map in \eqref{fr:two-slot-evaluation} has generic rank $R$.
\end{lemma}

\begin{proof}
Trivialize the two summands defining $Z$ near the selected base
points.  Let $z$ be the torus coordinate, whose value at $z_*$ is
$z_0\ne0$.  By the projective-bundle formula, the sections in
Lemma~\ref{fr:positive-characteristic-jets} are weight polynomials in $z$
with weight-$j$ coefficients in
\[
 H^0(W,B_p+jS)\otimes
 H^0(W,B_p+(pq-j)S),\qquad 0\le j\le pq.
\]
Let $A_*$ be the tensor product of the local algebras of the two
base Frobenius fibers.  The local quotient in
Lemma~\ref{fr:positive-characteristic-jets} is
\[
 A_*[z]/(z^p-z_0^p),
\]
free over $A_*$ with basis $1,z,\ldots,z^{p-1}$.  Project to the
coefficient of $1$ in this basis.  Exactly the weights $j=pa$
survive, with multipliers $z_0^{pa}$.  Their coefficients are the
values of the corresponding summands of
\eqref{fr:two-slot-evaluation}, in the chosen frames and the induced
twisted frames.  As the full jet map is surjective, these coefficients
span $A_*$, of dimension $p^{2n}=R$.  Thus the evaluation has full
rank at this pair of base points, and hence generically.  The
calculation uses a basis over the possibly nonreduced algebra $A_*$,
so reducedness of the Frobenius fibers is unnecessary.
\end{proof}

The jet calculation at $z_*$ proves generic full rank; the diagonal
estimate below will be tested at $(x',x')$, obtained from the
separately fixed point $x$.

\subsection{Sections on the Frobenius diagonal}
\label{fr:filtration}

The Frobenius fiber product has the cartesian square
\begin{equation}\label{fr:frobenius-square}
\begin{tikzcd}
 \Delta_F=W\times_{W'}W \arrow[r,"\mathrm{pr}_2"] \arrow[d,"\mathrm{pr}_1"'] & W \arrow[d,"F"]\\
 W \arrow[r,"F"'] & W'.
\end{tikzcd}
\end{equation}
Its reduced subscheme is the diagonal $W$.  Put
\[
 \mathcal L_p=
 \OO_{\Delta_F}(B_{p,1}+B_{p,2}+pqS_1).
\]
Write $h=F\circ\mathrm{pr}_1=F\circ\mathrm{pr}_2:
\Delta_F\to W'$. On this fiber product,
\[
 \OO(pS_1)\simeq h^*\OO_{W'}(S')\simeq\OO(pS_2).
\]
Consequently every pair of weights $a+b=q$ gives the same line bundle:
\[
 \OO_{\Delta_F}(B_{p,1}+B_{p,2}+paS_1+pbS_2)
 \simeq\mathcal L_p.
\]
On the diagonal of $W'\times W'$, every source twist of
\eqref{fr:two-slot-evaluation} becomes $-qS'$.
Finite base change and projection formula identify the twisted target as
\[
 (\mathcal E\otimes\mathcal E)(qS')\simeq h_*\mathcal L_p.
\]
The product sections defining the columns restrict to global sections
of this one bundle $\mathcal L_p$. Their values at a diagonal point
therefore lie in the image of $H^0(\Delta_F,\mathcal L_p)$ in the
corresponding fiber of $h_*\mathcal L_p$. This proves that the rank
at every diagonal point is at most
\begin{equation}\label{fr:diagonal-rank-upper}
 h^0(\Delta_F,\mathcal L_p).
\end{equation}

For a rank-$n$ bundle $V$ in characteristic $p$, denote by $A_j(V)$
the degree-$j$ part of its symmetric algebra modulo the $p$-th powers
of local generators.  This construction is independent of frame:
the $p$-th power of a linear combination is the sum of the $p$-th
powers in characteristic $p$.  Each $A_j(V)$ is locally free and a
quotient of $V^{\otimes j}$.  Set
\[
 u=n(p-1),\qquad e_j=\rk A_j(V),\qquad
 \sum_{j=0}^u e_j=p^n.
\]
Filtering $\mathcal L_p$ by powers of the ideal of the reduced
diagonal in $\Delta_F$ gives the graded bundles
\begin{equation}\label{fr:diagonal-graded}
 \mathcal G_j=\OO_W(2B_p+pqS)\otimes A_j(\Omega_W^1),
 \qquad 0\le j\le u.
\end{equation}
Indeed, in smooth local coordinates the algebra is generated by
parameter differences with their $p$-th powers zero, and its
degree-one conormal is $\Omega_W^1$.  This coordinate calculation
can be made étale locally or in completions and identifies the
global associated graded.  This is the diagonal-ideal form of the canonical Frobenius filtration
\cite[Corollary~3.5]{KS}. The same calculation, using
only the bundle from the second factor, filters $F^*\mathcal E$ with pieces
\cite[Theorem~3.7]{Sun}
\begin{equation}\label{fr:Frob-pullback-graded}
 \OO_W(B_p)\otimes A_j(\Omega_W^1).
\end{equation}

Complementary multiplication is a perfect pairing, giving
\begin{equation}\label{fr:socle-pairing}
 A_{u-j}(V)\simeq A_j(V)^\vee\otimes(\det V)^{p-1}.
\end{equation}
On monomials, each exponent vector pairs with its complement to
$(p-1,\ldots,p-1)$.  The top line transforms by the character
$(\det V)^{p-1}$: this is immediate on diagonal matrices and hence
 identifies the character of $\mathrm{GL}_n$ on that line.

\subsection{A uniform bound for the diagonal rank}
\label{fr:rank-bound}

We bound the sections of every graded bundle in
\eqref{fr:diagonal-graded} by its rank times the same scalar
polynomial in $p$. This uniformity lets us sum the ranks of the
graded pieces without an additional factor for their number.
We begin with a slope bound on the fixed curve; here slope means
ordinary degree divided by rank.

\begin{lemma}\label{fr:truncated-slope}
There is a constant $c\ge0$, independent of the sufficiently large
residual characteristic $p$ and of $j$, such that for
$V=\Omega_W^1|_C$ one has
\[
 \mu_{\max}(A_j(V))\le(p-1)K_W\cdot C+nc.
\]
\end{lemma}

\begin{proof}
Write $F_C$ for absolute Frobenius of $C$.  Langer's instability
estimate \cite[Corollary~2.5]{Langer} gives
\begin{equation}\label{fr:Langer-min}
 L_{\min}(V):=
 \lim_{e\to\infty}p^{-e}\mu_{\min}(F_C^{e*}V)
 \ge-\frac{c}{p-1}.
\end{equation}
To make the uniformity explicit, let $g$ be the fixed genus of $C$.
One can choose a nef line bundle $A_C$ of degree at most
$d_g=\max\{4g-2,0\}$ such that $T_C\otimes A_C$ is globally
generated; for $g\le1$, take $A_C=\OO_C$. Langer's bound is
\[
 \mu_{\min}(V)-L_{\min}(V)
 \le\frac{(\rk V-1)\deg A_C}{p-1}.
\]
Since $\rk V=n$ and $\mu_{\min}(V)\ge0$, we may take
$c=(n-1)d_g$, independently of the reduction.

We also need, for every integer $i\ge0$,
\begin{equation}\label{fr:tensor-min}
 \mu_{\min}(V^{\otimes i})\ge iL_{\min}(V).
\end{equation}
Fix $p$ and $i$.  Choose a line bundle $A_e$ of degree
$-\mu_{\min}(F_C^{e*}V)+O(1)$, rounded up with a fixed
genus-dependent additive constant, so that
$F_C^{e*}V\otimes A_e$ is globally generated.  This follows from
Serre duality: after subtracting any point, its minimum slope can
be made greater than $2g(C)-2$, which annihilates $H^1$ and makes
evaluation at that point surjective.  For every vector-bundle
quotient $Q$ of $V^{\otimes i}$, the bundle
$F_C^{e*}Q\otimes A_e^{\otimes i}$ is then globally generated.
Consequently
\[
 p^e\mu(Q)+i\deg A_e\ge0.
\]
Divide by $p^e$ and let $e\to\infty$ to prove
\eqref{fr:tensor-min}.  This limit is taken for each fixed $p,i$;
there is no interchange of Frobenius-iteration and characteristic
limits.

Since $A_{u-j}(V)$ is a quotient of $V^{\otimes(u-j)}$, the perfect
pairing \eqref{fr:socle-pairing} yields
\begin{align*}
 \mu_{\max}(A_j(V))
 &= (p-1)\deg\det V-\mu_{\min}(A_{u-j}(V))\\
 &\le(p-1)K_W\cdot C+\frac{(u-j)c}{p-1}\\
 &\le(p-1)K_W\cdot C+nc.
\end{align*}
\end{proof}

The divisors $B_p-pD=T_0-(p-k_0N_p)D$ range over a fixed finite
list. Combining~\eqref{fr:intersection-bounds},
\eqref{fr:flag-degrees}, and Lemma~\ref{fr:truncated-slope}
therefore gives, uniformly for $0\le j\le u$,
\begin{align}
 \mu_{\max}(\mathcal G_j|_C)
 &\le p(4+2/r)\Lambda H^n+O(1)\notag\\
 &\le M_p:=7p\Lambda H^n+c_0,\label{fr:common-slope-bound}
\end{align}
where $c_0\ge0$ is fixed.  The coefficient $7$ provides harmless
room in this common bound.

\begin{lemma}\label{fr:diagonal-sections}
For all sufficiently large $p$,
\[
 h^0(\Delta_F,\mathcal L_p)
 \le R\bigl(7^nH^n+O(1/p)\bigr)<R/4.
\]
The error constant is independent of the filtration index $j$.
\end{lemma}

\begin{proof}
Write the fixed flag as
\[
 W=W_n\supset W_{n-1}\supset\cdots\supset W_1=C,
 \qquad W_{n-k}\in|h_k|_{W_{n-k+1}}
 \quad(1\le k\le n-1).
\]
At step $k$, fix the preceding nonnegative integers
$b_1,\ldots,b_{k-1}$ and put
\[
 \mathcal F_k=
 \mathcal G_j|_{W_{n-k+1}}
 \Bigl(-\sum_{i<k}b_i h_i\Bigr).
\]
For each $b_k\ge0$, the divisor restriction sequence is
\[
 0\longrightarrow\mathcal F_k(-(b_k+1)h_k)
 \longrightarrow\mathcal F_k(-b_kh_k)
 \longrightarrow
 \mathcal F_k(-b_kh_k)|_{W_{n-k}}
 \longrightarrow0.
\]
Iterating it for $b_k=0,1,\ldots$ and then proceeding to the next
member of the flag bounds the section space by
\begin{equation}\label{fr:flag-lattice-sum}
 h^0(W,\mathcal G_j)\le
 \sum_{b_1,\ldots,b_{n-1}\ge0}
 h^0\left(C,\mathcal G_j|_C\Bigl(-\sum_i b_i h_i|_C\Bigr)\right).
\end{equation}
At each stage the remainder is zero after sufficiently negative
ample twisting.  Its stopping index need not be uniform in $j,p$
or in the preceding twists: enlarging the nonnegative finite sums
to the displayed lattice sum preserves the inequality.

A summand in \eqref{fr:flag-lattice-sum} vanishes when
$\sum_i b_i(h_i\cdot C)>M_p$, because then its maximum slope is
negative.  Each nonzero summand has at most $e_j(M_p+1)$ sections.
Indeed, evaluation at $\lfloor M_p\rfloor+1$ distinct points is injective;
the kernel has negative maximum slope after subtracting those
points.  By \eqref{fr:flag-degrees} it follows that
\begin{align}
 h^0(W,\mathcal G_j)
 &\le e_j(M_p+1)
       \prod_{i=1}^{n-1}
       \left(1+\frac{M_p}{l_i\Lambda H^n}\right)\notag\\
 &=e_jp^n\bigl(7^nH^n+O(1/p)\bigr).
 \label{fr:common-flag-polynomial}
\end{align}
The leading coefficient follows from
$\Lambda=\prod_i l_i$.  More importantly, the whole scalar
polynomial on the first line is the same for every $j$.

Sum over the filtration in \eqref{fr:diagonal-graded}.  The ranks
satisfy $\sum_j e_j=p^n$, so \eqref{fr:common-flag-polynomial} gives
the asserted bound with $R=p^{2n}$.  There is no extra factor for
the number of graded pieces.  Finally
\[
 7^nH^n\le(14\epsilon)^n<1/4
\]
by \eqref{fr:intersection-bounds}, proving the strict inequality
for all large $p$.
\end{proof}

\subsection{The two orders of the determinant}
\label{fr:determinant}

Choose $R$ generically independent columns of
\eqref{fr:two-slot-evaluation} from bases of its source vector spaces.
Their determinant is a nonzero section
\[
 0\ne\sigma\in H^0(W'\times W',\mathcal Q_1\boxtimes\mathcal Q_2)
\]
of an actual external-product line bundle. To specify its factors,
let $m_i$ be the sum of the weights of the $R$ chosen columns in
slot $i$. The target determinant is
$(\det\mathcal E)^s\boxtimes(\det\mathcal E)^s$, and the
chosen source determinant is
$\OO(-m_1S')\boxtimes\OO(-m_2S')$. Thus
\[
 \mathcal Q_i=(\det\mathcal E)^{\otimes s}
                 \otimes\OO_{W'}(m_iS'),\qquad
 \lambda_i=m_i/R.
\]
Every column weight lies between $0$ and $q$, so
$0\le\lambda_i\le q$ independently of the columns chosen.

Under the numerical identification with the base-field twist,
\begin{equation}\label{fr:slot-determinant-class}
 \frac{c_1(\mathcal Q_i)}R
 =\frac{c_1(\mathcal E)}s+\lambda_iS
 =\frac{B_p}p+\frac{p-1}{2p}K_W+\lambda_iS,
 \qquad 0\le\lambda_i\le q.
\end{equation}
The Frobenius first-Chern-class identity is also
\cite[Lemma~4.2]{Sun}; it is an identity of rational divisor classes.  For the second equality, use
\eqref{fr:Frob-pullback-graded}.  Complementary degrees in
\eqref{fr:socle-pairing} have total first Chern class
$e_j(p-1)K_W$.  Summing gives
\[
 \sum_jc_1(A_j(\Omega_W^1))=\frac{s(p-1)}2K_W.
\]
Thus $c_1(F^*\mathcal E)=sB_p+s(p-1)K_W/2$, and Frobenius pullback of
twisted divisor classes multiplies by $p$, as required.

For $0\le\lambda\le q$, the fixed movable inequality gives
\[
 (D+K_W/2+\lambda S)\cdot\gamma\le40\epsilon J\cdot\gamma,
\]
where pullback by $\pi_x$ is understood. This concerns only
intersections of fixed divisors and the fixed complete-intersection
witness, so it remains true on the reduction. The class in
\eqref{fr:slot-determinant-class} differs from its
left-hand divisor at $\lambda=\lambda_i$ by
\[
 \frac1p\bigl(B_p-pD-K_W/2\bigr).
\]
Its numerator belongs to the fixed finite list
$T_0-aD-K_W/2$, $a\in\{0,\ldots,k_0-1\}$. The resulting error after
intersection with $\gamma$ is therefore $O(p^{-1})$, uniformly
over the selected determinant columns.

Let $x'$ be the twisted point and let $0\ne\tau_i\in
H^0(W',\mathcal Q_i)$.  The strict transform of its effective divisor on
the blowup at $x'$ pairs nonnegatively with the twist of $\gamma$.
Since $J\cdot\gamma>0$, it follows that
\begin{equation}\label{fr:slot-order}
 \frac{\ord_{x'}(\tau_i)}R\le40\epsilon+O(1/p).
\end{equation}
The constant is uniform in the columns and in $\tau_i$.  This
argument tests sections on the reduction itself; it does not
require them to lift to characteristic zero.

K\"unneth identifies the space containing $\sigma$ with
$H^0(W',\mathcal Q_1)\otimes H^0(W',\mathcal Q_2)$.  Choose bases adapted to the
orders of vanishing at $x'$ in each factor.  In each degree their
leading terms are linearly independent.  Tensor products of these
leading terms remain independent in each bidegree of the two
disjoint sets of parameters; terms of different bidegrees cannot
cancel.  Every nonzero tensor consequently has order at most the
sum of the two maximum basis orders.  Applying
\eqref{fr:slot-order}, we obtain
\begin{equation}\label{fr:determinant-order-upper}
 \ord_{(x',x')}\sigma\le R\bigl(80\epsilon+O(1/p)\bigr).
\end{equation}

On the other hand, \eqref{fr:diagonal-rank-upper} and
Lemma~\ref{fr:diagonal-sections} bound the matrix rank at $(x',x')$
by a number strictly smaller
than $R/4$.  Trivialize its line bundles near that point.  Invertible
constant row operations make more than $3R/4$ rows vanish at the
point.  Each entry of those rows belongs to the maximal ideal, so
every term of the determinant has order at least $3R/4$.  Therefore
\begin{equation}\label{fr:determinant-order-lower}
 \ord_{(x',x')}\sigma\ge3R/4.
\end{equation}
Equations \eqref{fr:determinant-order-upper} and
\eqref{fr:determinant-order-lower} contradict
$80\epsilon<3/4$ for sufficiently large $p$.


This contradiction proves Theorem~\ref{thm:finite-data}.
The full-rank calculation used $z_*$, while the diagonal estimate
holds at every point and was tested at the separately chosen $x$.
The errors were controlled by one fixed finite list of divisor
classes and one fixed flag polynomial. Thus their constants do
not depend on $p$, the filtration index, or the chosen columns.

\subsection{Application to the terminal case}
\label{fr:application}

We return to Theorem~\ref{thm:terminal-big}. Thus $X$ is a projective
terminal $\Q$-factorial variety of dimension $n\ge2$ with
$K_X\sim0$, and $L$ is a nef Cartier divisor of full nef dimension.
Assume that $L$ is not big. Recall the ample rational divisors
$P=r(L+tA)$, where $A$ is a fixed ample Cartier divisor and $t>0$
tends to zero, with
\[
 r\longrightarrow\infty,\qquad P^n\longrightarrow\epsilon^n,
 \qquad \rho=(P^n)^{1/n}.
\]
The positive rational number $\epsilon$ was chosen so that
$(14\epsilon)^n<1/4$ and $80\epsilon<3/4$. Choose the integer $q$
supplied by Proposition~\ref{prop:large-seshadri}, and then choose
$t$ sufficiently small that this proposition and
Lemma~\ref{lemma:scalar-bound} apply, that $r>q+1$, and that
$P^n<(2\epsilon)^n$. We now fix $t$, and therefore also $P,r,\rho$
and the projective bundle in~\eqref{eq:bundle}. Every remaining
choice is made over $\C$ before reducing to positive characteristic.

\paragraph{The divisors and the flag.}
Choose a smooth projective resolution $w:W\to X$ and set
\[
 D=w^*P,\qquad S=w^*L,\qquad H_0=w^*P.
\]
The divisor $S$ is integral Cartier, and $H_0$ is nef and big.
Since $X$ is terminal and $K_X\sim0$, there is an integral
effective $w$-exceptional divisor $E$ with $K_W\sim E$. The
projection formula and the nefness of $P-rL$ give
\[
 K_WH_0^{n-1}=0,\qquad
 (2D+qS)H_0^{n-1}\le(2+q/r)H_0^n.
\]
At $H_0$, the inequalities
$K_WH_0^{n-1}<2H_0^n/r$ and
$(2D+qS)H_0^{n-1}<4H_0^n$ are strict. Thus, by continuity of
intersection products, a sufficiently small rational ample
perturbation $H=H_0+\eta A_W$, with $A_W$ an ample Cartier divisor
on $W$ and $\eta>0$, satisfies
\[
 H^n\le(2\epsilon)^n,\qquad
 K_WH^{n-1}\le2H^n/r,\qquad
 (2D+qS)H^{n-1}\le4H^n.
\]
This gives~\eqref{fr:intersection-bounds}.

The canonical divisor of $W$ is pseudo-effective. Generic
semipositivity of the cotangent bundle, originating in \cite{Miyaoka}
and used here in the form \cite[Theorem~2.1]{CP}, together with
restriction to sufficiently general complete-intersection curves
\cite[Theorem~6.1]{MR}, gives a smooth integral flag with
successive divisor classes $l_iH$, for sufficiently large divisible
integers $l_i$, such that
\[
 \mu_{\min}(\Omega_W^1|_C)\ge0.
\]
This is the flag in Hypothesis~\textup{(i)} of
Theorem~\ref{thm:finite-data}.

\paragraph{The finite jet system.}
On the bundle of~\eqref{eq:bundle}, write
$R=P_1+P_2+q\xi$. The divisor $R$ is ample: $\xi$ is nef and
relatively ample, and $P_1+P_2$ is the pullback of an ample
divisor on $X\times X$. Proposition~\ref{prop:large-seshadri} supplies
a smooth point $z_*$ in the complement of the two axes with
\[
 \varepsilon(R;z_*)>2n+2.
\]
The point can be chosen over the open set where $w$ is an
isomorphism. Choose a rational number $a$ strictly between
$2n+2$ and $\varepsilon(R;z_*)$. On the blowup
$b:\widetilde Z\to Z$ at $z_*$, with exceptional divisor $G$,
the rational divisor $b^*R-aG$ is ample. For sufficiently large
divisible $k_0$, Serre vanishing gives
\[
 H^1\bigl(\widetilde Z,
          \OO_{\widetilde Z}(k_0b^*R-k_0aG)\bigr)=0.
\]
Since $z_*$ is smooth, the sections on the thickened exceptional
divisor $k_0aG$ identify with the jet quotient at $z_*$ modulo
$\mathfrak m_{z_*}^{k_0a}$, in a local frame for $k_0R$.
The restriction sequence therefore shows that sections of
$k_0R$ generate jets through order $k_0a-1$, and hence through
order $(2n+2)k_0$. Increase the divisibility of $k_0$ to make
$k_0P$ Cartier. Pulling these sections to
\[
 Z_W=\PP\bigl(\OO_W(S)_1\oplus\OO_W(S)_2\bigr)
       \longrightarrow W\times W
\]
preserves the jets at the corresponding point, because the induced
map $Z_W\to Z$ is an isomorphism near that point. The pulled-back divisor is
$D_1+D_2+q\xi$. We have obtained the finite system required in
Hypothesis~\textup{(ii)}.

\paragraph{The movable curve class.}
Choose a very general point $x\in W$ outside the exceptional locus,
with image in the locus where Lemma~\ref{lemma:scalar-bound}
holds. Let $\pi_x:\widehat W\to W$ be its blowup and $J$ its
exceptional divisor. The rational divisor
\[
 D^+=2D+E
\]
is big. For every sufficiently divisible positive integer $m$,
normality of $X$ and exceptionality of $E$ give
$w_*\OO_W(mE)=\OO_X$. Thus every section of $mD^+$ comes from a
section of $2mP$, multiplied by the section of the fixed divisor
$mE$. Since $x\notin\Supp E$, Lemma~\ref{lemma:scalar-bound}
implies
\begin{equation}\label{fr:application-order}
 \ord_x(s)/m\le8\rho<16\epsilon
 \qquad\bigl(0\ne s\in H^0(W,\OO_W(mD^+))\bigr).
\end{equation}

It follows that
\begin{equation}\label{fr:application-not-pe}
 \pi_x^*D^+-40\epsilon J\quad\text{is not pseudo-effective}.
\end{equation}
Indeed, if it were pseudo-effective, a convex combination with
the big divisor $\pi_x^*D^+$ would make
$\pi_x^*D^+-cJ$ big for any rational $c$ with
$16\epsilon<c<40\epsilon$. A nonzero section of a sufficiently
divisible multiple would then give a section violating
\eqref{fr:application-order}.

By the duality between pseudo-effective divisors and strongly
movable curves~\cite{BDPP}, the strict separation
in~\eqref{fr:application-not-pe} is detected by a class
\[
 \gamma=f_*(A_1\cdots A_{n-1}),
 \qquad f:V\longrightarrow\widehat W,
\]
where $V$ is smooth projective, $f$ is birational, and the $A_i$
are ample integral Cartier divisors. To obtain this form, first
choose a generating strongly movable class with negative pairing,
then approximate its ample real classes by rational ones and
clear denominators. We have
\[
 (\pi_x^*D^+)\cdot\gamma<40\epsilon J\cdot\gamma.
\]
The left side is nonnegative, because $\pi_x^*D^+$ is
pseudo-effective, so $J\cdot\gamma>0$.

For every $0\le\lambda\le q$,
\[
 D^+-(D+K_W/2+\lambda S)
 \sim_{\R}w^*(P-\lambda L)+E/2
\]
is pseudo-effective: $P-\lambda L=(P-rL)+(r-\lambda)L$ is nef
and $E$ is effective. Pulling back to $\widehat W$ and pairing
with $\gamma$ therefore gives
\[
 \pi_x^*(D+K_W/2+\lambda S)\cdot\gamma
 \le(\pi_x^*D^+)\cdot\gamma
 <40\epsilon J\cdot\gamma.
\]
This proves Hypothesis~\textup{(iii)}. All the hypotheses of
Theorem~\ref{thm:finite-data} now hold for one fixed collection of
complex data, a contradiction. Hence $L$ is big, completing the
proof of Theorem~\ref{thm:terminal-big}.

\section{Completion of the full-dimension induction}
\label{sec:completion}

Theorem~\ref{thm:terminal-big} treats varieties with terminal
singularities and linearly trivial canonical divisor.  We now pass to
klt Calabi--Yau pairs and complete the full-dimension assertion.  The
boundary perturbation below follows the reduction in
\cite[Section~7]{LP}: an effective representative of a small boundary
perturbation plus the nef divisor allows ordinary log abundance to be
applied to a multiple of that nef divisor.

\begin{proposition}\label{prop:full-all}
Let $n\geq 2$, and assume $(P_j)$ for every $j<n$, together with
Theorems~\ref{input:abundance}
and~\ref{input:models}.  Let $(X,B)$ be a projective klt $\Q$-pair over
$\C$ of dimension $n$, and let $M$ be a nef $\Q$-Cartier divisor on $X$.
Suppose that $T=K_X+B$ is pseudo-effective and that, for some rational
$a>0$,
\[
  (T+aM)\cdot C>0
\]
for every integral curve $C$ through a very general point of $X$.
Then $T+cM$ is big for every rational $c>0$.
\end{proposition}

\begin{proof}
Proposition~\ref{prop:full-positive} proves the assertion when
$\kappa(T)>0$.  Proposition~\ref{prop:reduction-cy} reduces the remaining
case to the following claim in dimension $n$:
\begin{equation}\label{eq:completion-cy-claim}
 \begin{gathered}
 (X,B)\text{ projective klt},\qquad K_X+B\equiv 0,\\
 M\text{ nef and $\Q$-Cartier of full nef dimension}
 \quad\Longrightarrow\quad M\text{ big}.
 \end{gathered}
\end{equation}
We prove this claim first with $B=0$ and then with arbitrary boundary.
We may take a small projective $\Q$-factorialization throughout: the
adjoint and $M$ pull back, full nef dimension is preserved on the
isomorphism locus, and bigness of the pullback implies bigness downstairs.

\medskip
\noindent\emph{The case of zero boundary.}
Suppose that $X$ is klt and $K_X\equiv0$.  Theorem~\ref{input:abundance}
makes $K_X$ semiample, hence $K_X\sim_{\Q}0$.  Indeed, a semiample
numerically trivial divisor defines a morphism with zero-dimensional
image, so a sufficiently divisible multiple is linearly trivial.

Take the global index-one cover
\[
  \pi\colon X^{\mathrm{ind}}\longrightarrow X
\]
associated with a trivialization of the smallest positive multiple of
$K_X$ that is linearly trivial.  This is a finite quasi-\'etale cover,
and
\[
  K_{X^{\mathrm{ind}}}=\pi^*K_X\sim0.
\]
The cover is klt.  Since its canonical divisor is Cartier, its log
discrepancies are positive integers, hence at least one; it therefore
has canonical singularities.
These are the standard index-one cover properties for klt varieties;
see~\cite[Definition~5.19 and Proposition~5.20]{KM}.
By the crepant terminalization theorem
\cite[Corollary~1.4.3]{BCHM}, there is a projective birational morphism
\[
  h\colon \widetilde X\longrightarrow X^{\mathrm{ind}}
\]
with $\widetilde X$ terminal and $\Q$-factorial and
$K_{\widetilde X}=h^*K_{X^{\mathrm{ind}}}$.
In particular, $K_{\widetilde X}$ is Cartier and linearly trivial, as
required by Theorem~\ref{thm:terminal-big}.

Set $\widetilde M=(\pi\circ h)^*M$.  This divisor is nef and has full
nef dimension.  To verify the latter assertion, choose a very general
point upstairs outside the exceptional locus of $h$ and above the
very general locus used for $M$.  Every curve through this point has
a curve as its image in $X$, and the projection formula gives strictly
positive $\widetilde M$-degree.  A positive Cartier multiple of
$\widetilde M$ is therefore big by Theorem~\ref{thm:terminal-big}.
Bigness is detected by a dominant generically finite pullback, so $M$
is big on $X$.

\medskip
\noindent\emph{Introducing the boundary.}
Now let $(X,B)$ satisfy the hypotheses of
\eqref{eq:completion-cy-claim}, with $B\ne0$.
Theorem~\ref{input:abundance} gives
\[
  K_X+B\sim_{\Q}0.
\]
Choose a rational $\alpha>0$ sufficiently small that
$(X,(1+\alpha)B)$ remains klt, and put
\[
  T_*=K_X+(1+\alpha)B\sim_{\Q}\alpha B.
\]
Thus $T_*$ is pseudo-effective and $\kappa(T_*)\geq0$.  The full
condition holds for $T_*$ and $M$: for a curve $C$ through a point
outside $\Supp B$, we have $T_*\cdot C=\alpha B\cdot C\geq0$,
whereas $M\cdot C>0$ if that point is very general.
We will show that $T_*+M$ is big without using the conclusion
of the present proposition.

If $\kappa(T_*)>0$, this follows from
Proposition~\ref{prop:full-positive}.  Suppose instead that
$\kappa(T_*)=0$, and apply
Lemma~\ref{lemma:two-nef-model} to $(X,(1+\alpha)B)$ and $M$.
Write $Y$ for its final model, $B_Y$ and $N$ for the transforms of
$B$ and $M$, and
\[
  U=K_Y+(1+\alpha)B_Y.
\]
The output has $U\sim_{\Q}0$ and, for some rational $s>0$, the divisor
$U+sN$ is nef of full nef dimension.  The maps in the construction
extract no divisors, so pushing forward the rational linear equivalence
$K_X+B\sim_{\Q}0$ also gives
\[
  K_Y+B_Y\sim_{\Q}0.
\]
Subtracting these two equivalences shows that
$\alpha B_Y\sim_{\Q}0$.  Since $B_Y$ is effective on a projective
variety, intersection with a sufficiently ample complete-intersection
curve forces $B_Y=0$.  The ordinary pair on $Y$ is therefore klt with
zero boundary and $K_Y\sim_{\Q}0$.
Moreover, $N$ is nef of full nef dimension, because
$U+sN\equiv sN$.  The zero-boundary case just proved makes $N$, and
hence $U+sN$, big.

On a common smooth resolution, with maps $p$ to $X$ and $q$ to $Y$,
the effective exceptional comparison from the reduction is
\[
  p^*(T_*+sM)=q^*(U+sN)+E,
  \qquad E\geq0.
\]
It follows that $T_*+sM$ is big.  Choose a positive rational $\lambda$
with $\lambda<\min\{1,1/s\}$.  Then
\[
  T_*+M
   =\lambda(T_*+sM)+(1-\lambda)T_*+(1-\lambda s)M
\]
is big, since the first summand is big and the other two are
pseudo-effective.  We have thus established bigness of $T_*+M$ in
both cases.

Choose an effective $\Q$-divisor $G$ with
\[
  G\sim_{\Q}T_*+M\sim_{\Q}\alpha B+M.
\]
For a sufficiently small rational $u>0$, the divisor
\[
  \Delta_u=(1-u\alpha)B+uG
\]
is effective and $(X,\Delta_u)$ is klt.  To see this explicitly, take
$u\alpha<1$ and choose $u$ small enough that $(X,B+uG)$ is klt on a
fixed log resolution; then $\Delta_u\leq B+uG$.
Its adjoint satisfies
\[
  K_X+\Delta_u
   \sim_{\Q} K_X+B+u(G-\alpha B)
   \sim_{\Q}uM.
\]
This adjoint is nef, so Theorem~\ref{input:abundance} makes it
semiample.  Consequently $M$ itself is semiample.
The morphism given by a free multiple of $M$ has zero-dimensional
general fiber: a positive-dimensional fiber through a very general
point would contain a curve of $M$-degree zero, contrary to full nef
dimension.  That morphism is therefore generically finite onto its
image, and $M$ is big.  This proves
\eqref{eq:completion-cy-claim} for all klt Calabi--Yau pairs.

Proposition~\ref{prop:reduction-cy} now supplies the case
$\kappa(T)=0$, while Proposition~\ref{prop:full-positive} supplies
$\kappa(T)>0$.  These exhaust the possibilities furnished by the
minimal-model and abundance inputs for a pseudo-effective klt
adjoint, and establish the proposition.
\end{proof}

\section{Positive parts and numerical semiampleness}
\label{sec:positive}

We complete the simultaneous induction by proving the following
positive-part statement.

\begin{proposition}\label{prop:positive-part}
For every projective klt $\Q$-pair $(X,B)$ over $\C$ with
$K_X+B$ pseudo-effective and every nef $\Q$-Cartier divisor $M$ on
$X$, there is a smooth projective birational model $w\colon W\to X$
such that
\[
                     P_\sigma\bigl(w^*(K_X+B+M)\bigr)
\]
is rational and numerically equivalent to a semiample $\Q$-divisor.
\end{proposition}

For the inductive step, fix such $(X,B)$ and $M$ with $\dim X=n\geq2$,
and put $T=K_X+B$.  We assume $(Z_j)$ for $j<n$ and use $(P_n)$
from Proposition~\ref{prop:full-all}.  Following the positive-part
method of Lazi\'c--Peternell~\cite[Theorem 5.3]{LP}, we first make
a sufficiently large adjoint nef.  Full nef dimension then gives
bigness by $(P_n)$; otherwise its nef reduction permits the nef
summand to descend to a lower-dimensional base.

\subsection{Positive parts under the required comparisons}

We will repeatedly use the following properties of Nakayama's
divisorial Zariski decomposition \cite{Nakayama}.

\begin{lemma}\label{lemma:positive-properties}
The following assertions hold.
\begin{enumerate}
\item Let $p\colon V\to X$ be a projective birational morphism from
a smooth projective variety to a normal projective variety.  If $D$
is a pseudo-effective $\Q$-Cartier divisor on $X$ and $E$ is an
effective $p$-exceptional rational divisor, then
\[
                       P_\sigma(p^*D+E)=P_\sigma(p^*D).
\]
\item Let $f\colon V\to S$ be a surjective morphism of smooth
projective varieties, and let $D$ be a pseudo-effective rational
divisor on $S$.  If $P_\sigma(D)$ is nef, then
\[
                         P_\sigma(f^*D)=f^*P_\sigma(D).
\]
\item For rational divisors on a fixed smooth projective variety,
rationality and numerical semiampleness of the positive part depend
only on the numerical class of the divisor.
\end{enumerate}
\end{lemma}

\begin{proof}
The first two statements are \cite[Lemmas 2.4 and 2.5]{LP}; the first
allows the target $X$ to be normal.  For the third, each asymptotic
divisorial order $\sigma_\Gamma(D)$ depends only on the numerical class.
Thus numerically equivalent divisors have the same negative part as
an actual real divisor.  For rational $D,D'$ with $D\equiv D'$,
\[
                  P_\sigma(D')-P_\sigma(D)=D'-D.
\]
Consequently one positive part is rational if and only if the other
is, and the two positive parts are numerically equivalent.
\end{proof}

In particular, a positive part that is rational and numerically
semiample is nef, so these properties persist on higher smooth
models.  Suppose that on a common smooth resolution of a birational
contraction one has
\begin{equation}\label{eq:positive-transfer}
                         p^*D\equiv q^*D'+E,
\end{equation}
where $D,D'$ are pseudo-effective rational Cartier divisors and $E$
is effective and $q$-exceptional.  If a smooth model over the target
has rational numerically semiample positive part for $D'$, take a
higher common resolution dominating that model.  The second part of
Lemma~\ref{lemma:positive-properties} pulls that positive part to the
common resolution.  Its first and third parts, applied to
\eqref{eq:positive-transfer}, give the same conclusion for $D$.
This is the comparison that will be used at each birational step.

\subsection{An MMP trivial on the nef divisor}

\begin{lemma}\label{lemma:large-nef-adjoint}
Let $(X,B)$ be a projective $\Q$-factorial klt $\Q$-pair of
dimension $n$, with $T=K_X+B$ pseudo-effective, and let $L$ be a nef
Cartier divisor on $X$.  If $m$ is an integer greater than $2n$,
there is a terminating $(T+mL)$-MMP whose every step is $L$-trivial
and is an ordinary $T$-negative step.  On its output, the transform
of $L$ is Cartier and nef, and the transform of $T+mL$ is nef.
\end{lemma}

\begin{proof}
Use Theorem~\ref{input:models} for the generalized klt pair whose
nef data are $mL$ on $X$.  Its adjoint $T+mL$ is pseudo-effective.
Suppose inductively that on the current model the ordinary pair is
klt and the transform, still denoted by $L$, is Cartier and nef.
Every $(T+mL)$-negative extremal ray is $T$-negative.  By the ordinary
length bound, it has a rational curve generator $C$ with
$0<-T\cdot C\leq2n$.  If $L\cdot C>0$, integrality gives
\[
                      (T+mL)\cdot C\geq-2n+m>0,
\]
which is impossible.  Hence $L\cdot C=0$.

The Cartier descent statement in the ordinary cone and contraction
theorem gives a Cartier divisor on the contraction base whose
pullback is $L$.  It is nef: lift any curve on the base to a curve
dominating it and use the projection formula.  In a flip its
pullback to the flipped variety is therefore again Cartier and nef.
The relative signs of $T+mL$ are those of $T$, so the flip is also
an ordinary flip for the klt pair.  Thus the inductive hypotheses
persist.  The terminating program supplied by
Theorem~\ref{input:models} has all the stated properties.
\end{proof}

Apply the lemma after a small $\Q$-factorialization, with
$L=m_0M$ for an integer $m_0>0$ making $L$ Cartier.  The program is
$M$-trivial, so on a common resolution its ordinary adjoint
comparison also gives
\[
                     p^*(T+M)=q^*(T'+M')+E,
\]
where $E$ is effective and exceptional over the output.  The
transformed ordinary adjoint $T'$ remains pseudo-effective, and
$M'$ is nef.  By \eqref{eq:positive-transfer}, it suffices to prove
$(Z_n)$ on that output.  We may consequently retain the notation
$X,B,T,M,L$ and assume in addition that
\begin{equation}\label{eq:large-adjoint-nef}
                               T+mL\text{ is nef}.
\end{equation}

\subsection{The case of full nef dimension}

Suppose first that $T+mL$ has full nef dimension.  Then
\eqref{eq:full-condition} holds, and $(P_n)$ shows that $T+M$ is big.
Write, by Kodaira's lemma,
\[
                     T+M\sim_\Q A_1+G_1,
\]
where $A_1$ is ample and $G_1\geq0$ is rational.  Choose a rational
$v>0$ sufficiently small that $(X,B+vG_1)$ is klt.  Since
$M+vA_1$ is ample, a general effective rational representative
$A'$ of this divisor can be chosen with sufficiently small
coefficients that $(X,B+vG_1+A')$ is klt.  Its adjoint satisfies
\begin{equation}\label{eq:big-ordinary-boundary}
 K_X+B+vG_1+A'\sim_\Q(1+v)(T+M).
\end{equation}
It is big and has a good minimal model: one may use the big-case
theorem of \cite{BCHM}, or combine
Theorems~\ref{input:models} and \ref{input:abundance}.
On a common resolution, the left side of
\eqref{eq:big-ordinary-boundary} is the pullback of a semiample
divisor plus an effective exceptional divisor.  Its positive part
is therefore the semiample pullback.  Dividing by $1+v$ proves
$(Z_n)$ in this case.

\subsection{Descent along a nef reduction}

We treat the remaining case, where the nef dimension of $T+mL$ is
strictly smaller than $n$.  Let
\[
                         X\dashrightarrow Z
\]
be its nef reduction, with $Z$ normal and projective.  A very general
compact fibre $F$ is contained in the morphism locus and has positive
dimension.  On this fibre,
\[
                         (T+mL)|_F\equiv0.
\]
The restriction $T|_F$ is pseudo-effective.  To see this, choose
effective rational approximations to $T$ after adding positive ample
perturbations tending to zero, and choose $F$ outside the countably
many exceptional choices for these approximations.  Their
restrictions are effective, and their classes approach $T|_F$.
After a resolution of $F$ the same argument applies to the pullbacks.
As $L|_F$ is nef and $T|_F\equiv-mL|_F$, both divisors are
numerically trivial.  Indeed, intersect this equality with an ample
complete-intersection curve on a smooth resolution of $F$.
Pseudo-effectivity makes the left side nonnegative and nefness makes
the right side nonpositive.  Their common value is zero, and a nef
divisor with zero intersection against an ample complete-intersection
curve is numerically trivial.

Restriction to a sufficiently general fibre gives a klt pair
$(F,B|_F)$ with adjoint $T|_F$.  By
Theorem~\ref{input:abundance}, its numerically trivial adjoint is
semiample, hence
\begin{equation}\label{eq:fibre-adjoint-trivial}
                         K_F+B|_F\sim_\Q0.
\end{equation}
These restrictions are made where the almost holomorphic nef
reduction is a morphism.  Normality and the klt restriction statement
hold after shrinking the base; alternatively they follow by
restricting a log resolution over a general smooth open.  Numerical
triviality, initially obtained for very general fibres, also follows
on a general fibre by the constancy of the relevant intersection
numbers in this family.

We next use the birational descent of nef divisors
\cite[Lemma 3.1]{LP}.  After resolving the nef reduction and modifying
its base birationally, it supplies a diagram
\[
                       X\xleftarrow{p}X_1\xrightarrow{f}Z_1
\]
in which $X_1,Z_1$ are smooth and projective, $p$ is birational,
$f$ is a contraction, and there is a rational divisor $M_1$ on
$Z_1$ with
\begin{equation}\label{eq:nef-descent}
                              p^*M\equiv f^*M_1.
\end{equation}
The divisor $M_1$ is nef, since its pullback is numerically
equivalent to the nef divisor $p^*M$ and nefness descends under a
surjective morphism.  We may take a further log resolution of the
pair on top without affecting \eqref{eq:nef-descent}.

Write the crepant boundary as $B_1^+-B_1^-$, with effective rational
divisors of disjoint support, so that
\[
                   K_{X_1}+B_1^+-B_1^-=p^*(K_X+B).
\]
Set $B_1=B_1^+$ and $E=B_1^-$.  Since the original boundary is
effective, $E$ is $p$-exceptional.  Klt singularities make all
coefficients of $B_1$ less than one, and its support is simple normal
crossings.  Thus $(X_1,B_1)$ is klt and
\begin{equation}\label{eq:resolved-boundary}
                    K_{X_1}+B_1=p^*T+E,\qquad E\geq0.
\end{equation}
In particular its adjoint is pseudo-effective.

If $F_1$ is a general fibre of $f$, it maps birationally onto the
corresponding compact fibre $F$ of the nef reduction.  Every
component of $E$ that meets $F_1$ restricts to an exceptional divisor
over $F$: for a centre dominating the base, codimension at least two
is preserved on the general fibre, and the other centres are avoided.
Equations~\eqref{eq:fibre-adjoint-trivial} and
\eqref{eq:resolved-boundary} therefore show that
\begin{equation}\label{eq:fibre-kappa-zero}
               \kappa\bigl(F_1,(K_{X_1}+B_1)|_{F_1}\bigr)=0.
\end{equation}
The restriction is an effective exceptional divisor up to rational
linear equivalence.  Equivalently, its section spaces are those of
the trivial divisor on the normal fibre $F$.  In particular the
general fibre pair has pseudo-effective adjoint and has a good
minimal model by Theorems~\ref{input:models} and
\ref{input:abundance}.

\subsection{A relative good model and a lower-dimensional adjoint}

At this point the nef summand comes from $Z_1$, while the ordinary
adjoint has Iitaka dimension zero on its general fibres.  We use a
relative good model to express that ordinary adjoint as a pullback
from a variety of dimension less than $n$.

The relative log canonical ring of the rational klt pair
$(X_1,B_1)$ over $Z_1$ is finitely generated.  More precisely, for
a positive integer $r$ making $r(K_{X_1}+B_1)$ Cartier, its Veronese
algebra
\[
 \bigoplus_{\ell\geq0}
 f_*\OO_{X_1}\bigl(\ell r(K_{X_1}+B_1)\bigr)
\]
is a finitely generated $\OO_{Z_1}$-algebra.  This is the relative
finite-generation theorem for rational klt pairs, in the form
\cite[Theorem 1.4]{FujinoFG}, obtained from the finite-generation theorem
of \cite{BCHM} and the canonical-bundle reduction of
Fujino--Mori \cite[Theorem~5.2]{FM}.
Together with the good minimal models of the very general closed
fibres established above, this verifies both hypotheses of
\cite[Theorem 2.12]{HX}.  That theorem gives a good minimal model
over $Z_1$; the termination statement for an MMP with scaling once
a good model exists, \cite[Corollary 2.9]{HX}, supplies a relative
program
\[
                       (X_1,B_1)\dashrightarrow(X_2,B_2).
\]
Thus $X_2$ is projective and $\Q$-factorial, $(X_2,B_2)$ is klt,
and $K_{X_2}+B_2$ is semiample over $Z_1$.  This application uses
very general closed fibres over $\C$ and requires no change of the
ground field to a geometric generic field.

By \eqref{eq:nef-descent}, we may replace the nef term numerically
by $f^*M_1$ before making the positive-part comparisons.  This
pullback is trivial on every step of the relative program and is
carried unchanged from the base.  The ordinary adjoint remains
pseudo-effective by numerical pushforward.

Let
\[
                    X_2\xrightarrow{g}Z_2\xrightarrow{j}Z_1
\]
be the relative semiample fibration of $K_{X_2}+B_2$, with $g$ a
contraction.  Its construction gives a $j$-ample $\Q$-Cartier
divisor $A_2$ on $Z_2$ such that
\begin{equation}\label{eq:relative-iitaka}
                         K_{X_2}+B_2\sim_\Q g^*A_2.
\end{equation}
The relative program preserves Iitaka dimension on the general
fibres, by its effective exceptional pullback comparisons.  From
\eqref{eq:fibre-kappa-zero}, the relative Iitaka dimension is zero.
Hence
\[
                           \dim Z_2=\dim Z_1<n.
\]

Apply Ambro's canonical bundle formula consequence
\cite[Theorem 0.2]{Ambro} to \eqref{eq:relative-iitaka}.  Its
hypotheses hold: $(X_2,B_2)$ is a projective klt pair with effective
rational boundary, $g$ is a contraction to a normal projective
variety, and the ordinary adjoint is rationally linearly pulled back
from a $\Q$-Cartier divisor.  It gives an effective rational boundary
$B_{Z_2}$ such that $(Z_2,B_{Z_2})$ is klt and
\begin{equation}\label{eq:base-adjoint}
              K_{X_2}+B_2\sim_\Q g^*(K_{Z_2}+B_{Z_2}).
\end{equation}
The base adjoint is pseudo-effective.  One direct verification,
using only the inputs already in force, is as follows.  The
pseudo-effective ordinary adjoint on $X_2$ has a good minimal model
by Theorems~\ref{input:models} and \ref{input:abundance}; it therefore
has a nonzero section in a sufficiently divisible positive multiple.
Since $g_*\OO_{X_2}=\OO_{Z_2}$, the projection formula and
\eqref{eq:base-adjoint} identify this section space with the
corresponding section space on $Z_2$.  Thus $K_{Z_2}+B_{Z_2}$ even
has an effective rationally linearly equivalent divisor.

We are now in the setting of $(Z_{\dim Z_2})$: the pair on the base
is klt with pseudo-effective adjoint, and $j^*M_1$ is a nef rational
Cartier divisor.  There is consequently a smooth projective
birational model $w_2\colon W_2\to Z_2$ on which
\[
                 P_\sigma\bigl(w_2^*(K_{Z_2}+B_{Z_2}+j^*M_1)\bigr)
\]
is rational and numerically semiample.  Resolve the main component
of $X_2\times_{Z_2}W_2$ to obtain a smooth projective variety $V$
mapping birationally to $X_2$ and surjectively to $W_2$.  Since the
displayed positive part is nef, the pullback formula of
Lemma~\ref{lemma:positive-properties}, together with
\eqref{eq:base-adjoint}, proves the same positive-part assertion on
$V$ for
\[
                         K_{X_2}+B_2+g^*j^*M_1.
\]

Take a common resolution dominating both $V$ and $X_1$.  The
relative MMP compares $K_{X_1}+B_1$ with $K_{X_2}+B_2$ by adding
an effective divisor exceptional over $X_2$.  The $M_1$ terms have
equal pullbacks, since they come from the relative base.  The
comparison \eqref{eq:positive-transfer} therefore proves the
positive-part assertion on a smooth model over $X_1$ for
$K_{X_1}+B_1+f^*M_1$.  Finally,
\eqref{eq:nef-descent} and \eqref{eq:resolved-boundary} give
\[
                 K_{X_1}+B_1+f^*M_1\equiv p^*(T+M)+E,
\]
where $E$ is effective and $p$-exceptional.  Another application of
Lemma~\ref{lemma:positive-properties} proves $(Z_n)$ on a smooth
model over the original $X$.

\begin{proof}[Completion of the proof of Proposition~\ref{prop:positive-part}]
The preceding argument proves $(Z_n)$ from $(P_n)$ and the
lower-dimensional assertions.  Together with
Proposition~\ref{prop:full-all} and the dimension-zero and
dimension-one cases, it completes the simultaneous induction.
\end{proof}

\subsection{Descent to the original variety}

The positive-part statement is birational, whereas the main theorem
requires a semiample representative on the original variety.  The
following descent uses the rational singularities of klt pairs;
it is the numerical semiampleness statement of
\cite[Lemma 2.14]{LP}.

\begin{lemma}\label{lemma:numsemiample-descent}
Let $p\colon W\to X$ be a projective birational morphism from a
smooth projective variety to a normal projective variety with
rational singularities, over $\C$.  Let $D$ be a $\Q$-Cartier
divisor on $X$.  If $p^*D$ is numerically equivalent to a semiample
$\Q$-divisor on $W$, then $D$ is numerically equivalent to a
semiample $\Q$-Cartier divisor on $X$.
\end{lemma}

\begin{proof}
Choose a semiample rational divisor $A_W$ with $A_W\equiv p^*D$.
After clearing denominators, the difference of the corresponding
line bundles is numerically trivial.  Numerically trivial line
bundles modulo $\Pic^0(W)$ are torsion.  Thus, after increasing a
positive integer $r$, we have
\[
                 \OO_W(rA_W)\simeq p^*\OO_X(rD)\otimes P_W
                 \quad\text{with }P_W\in\Pic^0(W).
\]
Rational singularities identify $\Pic^0(X)$ with $\Pic^0(W)$ by
pullback.  Write $P_W=p^*P_X$ and set
$\mathcal A_X=\OO_X(rD)\otimes P_X$.  Then
$p^*\mathcal A_X\simeq\OO_W(rA_W)$.

A positive power of the line bundle on the right is globally
generated.  Since $p_*\OO_W=\OO_X$, the projection formula
identifies its sections with those of the same power of
$\mathcal A_X$.  These sections generate downstairs: a failure at
any point of $X$ would remain a failure at every point above it.
Hence $\mathcal A_X$ is semiample.  A Cartier divisor representing
$\mathcal A_X$, divided by $r$, is semiample and numerically
equivalent to $D$, because $P_X$ is algebraically trivial.
\end{proof}

We can now prove Theorem~\ref{thm:main} over $\C$.  Its divisor
$D=K_X+B+M$ is nef, so
\[
                         P_\sigma(w^*D)=w^*D
\]
on every smooth birational model.  Proposition~\ref{prop:positive-part}
makes this pullback numerically semiample.  A klt variety has rational
singularities, so Lemma~\ref{lemma:numsemiample-descent} gives the
required semiample rational Cartier divisor on $X$.

\subsection{Algebraically closed fields of characteristic zero}

We finish by transferring the result to the ground fields in the
statement of Theorem~\ref{thm:main}.  The very general points used in
the proof were needed only over $\C$; numerical semiampleness itself
descends by a finite-type parameter argument.

\begin{lemma}\label{lemma:field-transfer}
Let $k_0\subset K$ be an extension of algebraically closed fields
of characteristic zero.  Let $X_0$ be a normal projective variety
over $k_0$, and let $D_0$ be a $\Q$-Cartier divisor.  If its pullback
$D_K$ to $X_K$ is numerically equivalent to a semiample
$\Q$-Cartier divisor, then $D_0$ has that property over $k_0$.
\end{lemma}

\begin{proof}
Clear denominators in $D_0$ and in a semiample representative over
$K$, and take a further multiple which is globally generated.  Its difference
from the corresponding multiple of $D_K$ is numerically trivial.
After another positive multiple this difference belongs to
$\Pic^0(X_K)$, since numerical triviality modulo $\Pic^0$ is
torsion.  We have therefore fixed an integer $r>0$ and a line
bundle $P_K\in\Pic^0(X_K)$ such that
\[
                       \OO_{X_K}(rD_K)\otimes P_K
\]
is globally generated.

Choose a $k_0$-point of $X_0$ and rigidify line bundles there.  The
finite-type scheme $\Pic^0(X_0)$ then has a Poincar\'e family.  Tensor
this family by the fixed line bundle $\OO_{X_0}(rD_0)$.  The set of
parameters for which the resulting line bundle is globally
generated is constructible.  Indeed, stratify the parameter scheme
so that the zeroth direct image is locally free and cohomology and
base change hold.  On each stratum use the evaluation map; the
image in the parameter space of the support of its cokernel is
closed, because projection from $X_0$ is proper.  Its complement
is the globally generated locus on that stratum.

Formation of the Picard scheme and of this constructible condition
commutes with the field extension.  The locus is nonempty after
extension to $K$, by the chosen $P_K$, and is consequently nonempty
over $k_0$.  As $k_0$ is algebraically closed, it has a $k_0$-point.
The corresponding globally generated line bundle, represented by
a Cartier divisor and divided by $r$, is the desired semiample
rational divisor numerically equivalent to $D_0$.
\end{proof}

\begin{proof}[Completion of the proof of Theorem~\ref{thm:main}]
Let $k$ be the given algebraically closed field of characteristic
zero.  Descend $X,B,M$, the required rational Cartier structures,
and a projective embedding to a finitely generated subfield of $k$.
Let $k_0\subset k$ be its algebraic closure inside $k$, and choose
an embedding $k_0\hookrightarrow\C$.

The hypotheses persist under the resulting algebraically closed
field extensions.  For klt singularities this follows from a log
resolution, which can be included among the descended data.
Nefness is invariant under field extension, since it can be tested
by ampleness after arbitrarily small rational ample perturbations.
Pseudo-effectivity is invariant for the same reason with bigness
in place of ampleness: a rational divisor is pseudo-effective
exactly when all its positive rational ample perturbations are
big, and bigness is preserved and reflected by extension of the
section spaces.  These tests show both that the descended data
satisfy the hypotheses and that their complex extension does.

The result over $\C$ supplies a semiample numerical representative
of $K_X+B+M$ after this extension.  Lemma~\ref{lemma:field-transfer}
descends it to $k_0$.  Extending that representative to $k$ preserves
global generation of a suitable multiple and numerical
equivalence.  It is the required semiample $\Q$-Cartier divisor
on the original variety.
\end{proof}

\begin{thebibliography}{99}

\bibitem{Ambro}
F.~Ambro,
\emph{The moduli $b$-divisor of an lc-trivial fibration},
Compos. Math. \textbf{141} (2005), no.~2, 385--403.
\href{https://doi.org/10.1112/S0010437X04001071}{DOI}.

\bibitem{NefReduction}
Th.~Bauer, F.~Campana, Th.~Eckl, S.~Kebekus, Th.~Peternell,
S.~Rams, T.~Szemberg, and L.~Wotzlaw,
\emph{A reduction map for nef line bundles},
in \emph{Complex Geometry: Collection of Papers Dedicated to Hans Grauert},
Springer, Berlin, 2002, pp.~27--36.
\href{https://doi.org/10.1007/978-3-642-56202-0_2}{DOI}.

\bibitem{BCHM}
C.~Birkar, P.~Cascini, C.~D.~Hacon, and J.~McKernan,
\emph{Existence of minimal models for varieties of log general type},
J. Amer. Math. Soc. \textbf{23} (2010), no.~2, 405--468.
\href{https://doi.org/10.1090/S0894-0347-09-00649-3}{DOI}.

\bibitem{BZ}
C.~Birkar and D.-Q.~Zhang,
\emph{Effectivity of Iitaka fibrations and pluricanonical systems of polarized pairs},
Publ. Math. Inst. Hautes \'Etudes Sci. \textbf{123} (2016), 283--331.
\href{https://doi.org/10.1007/s10240-016-0080-x}{DOI}.

\bibitem{Blanc}
J.~Blanc,
\emph{Algebraic structures of groups of birational transformations},
in \emph{Algebraic Groups: Structure and Actions} (M.~Bilen Can, ed.),
Proc.~Sympos.~Pure Math. \textbf{94}, American Mathematical Society,
Providence, RI, 2017, pp.~17--30.
\href{https://algebra.dmi.unibas.ch/blanc/articles/structurebir.pdf}{Author-hosted published offprint}.

\bibitem{BDPP}
S.~Boucksom, J.-P.~Demailly, M.~P\u{a}un, and Th.~Peternell,
\emph{The pseudo-effective cone of a compact K\"ahler manifold and varieties of negative Kodaira dimension},
J.~Algebraic Geom. \textbf{22} (2013), no.~2, 201--248.
\href{https://doi.org/10.1090/S1056-3911-2012-00574-8}{DOI}.

\bibitem{CP}
F.~Campana and M.~P\u{a}un,
\emph{Orbifold generic semi-positivity: an application to families of canonically polarized manifolds},
Ann.~Inst.~Fourier (Grenoble) \textbf{65} (2015), no.~2, 835--861.
\href{https://doi.org/10.5802/aif.2945}{DOI}.

\bibitem{Chaudhuri}
Priyankur Chaudhuri,
\emph{An inductive approach to generalized abundance using nef reduction},
C. R. Math. Acad. Sci. Paris \textbf{361} (2023), 417--421.
\href{https://doi.org/10.5802/crmath.420}{DOI}.

\bibitem{EKL}
L.~Ein, O.~K\"uchle, and R.~Lazarsfeld,
\emph{Local positivity of ample line bundles},
J. Differential Geom. \textbf{42} (1995), no.~2, 193--219.
\href{https://doi.org/10.4310/jdg/1214457231}{DOI}.

\bibitem{Fontanari}
Claudio Fontanari,
\emph{A remark on generalized abundance for surfaces},
Eur. J. Math. \textbf{10} (2024), article~7.
\href{https://doi.org/10.1007/s40879-023-00716-y}{DOI}.

\bibitem{FujinoFG}
O.~Fujino,
\emph{Birkar--Cascini--Hacon--McKernan},
author notes, version~1.02, 28 November 2009, Theorem~1.4.
\href{https://www.math.kyoto-u.ac.jp/~fujino/BCHM01.pdf}{Author's notes}.

\bibitem{FG}
O.~Fujino and Y.~Gongyo,
\emph{On the moduli $b$-divisors of lc-trivial fibrations},
Ann. Inst. Fourier (Grenoble) \textbf{64} (2014), no.~4, 1721--1735.
\href{https://doi.org/10.5802/aif.2894}{DOI}.

\bibitem{FH}
O.~Fujino and K.~Hashizume,
\emph{Existence of log canonical modifications and its applications},
Eur. J. Math. \textbf{9} (2023), no.~1, Paper No.~13, 32~pp.
\href{https://doi.org/10.1007/s40879-023-00598-0}{DOI}.

\bibitem{FM}
O.~Fujino and S.~Mori,
\emph{A canonical bundle formula},
J. Differential Geom. \textbf{56} (2000), no.~1, 167--188.

\bibitem{Fukuda}
S.~Fukuda,
\emph{An elementary semi-ampleness result for log canonical divisors},
Hokkaido Math. J. \textbf{40} (2011), no.~3, 357--360.
\href{https://doi.org/10.14492/hokmj/1319595860}{DOI}.

\bibitem{HX}
C.~D.~Hacon and C.~Xu,
\emph{Existence of log canonical closures},
Invent. Math. \textbf{192} (2013), no.~1, 161--195.
\href{https://doi.org/10.1007/s00222-012-0409-0}{DOI}.
Numbering follows arXiv:1105.1169v3.

\bibitem{Hanamura}
M.~Hanamura,
\emph{Structure of birational automorphism groups, I: non-uniruled varieties},
Invent.~Math. \textbf{93} (1988), no.~2, 383--403.
\href{https://doi.org/10.1007/BF01394338}{DOI}.

\bibitem{KS}
Y.~Kitadai and H.~Sumihiro,
\emph{Canonical filtrations and stability of direct images by Frobenius morphisms},
Tohoku Math. J. (2) \textbf{60} (2008), no.~2, 287--301.
\href{https://doi.org/10.2748/tmj/1215442876}{DOI}.

\bibitem{KM}
J.~Koll\'ar and S.~Mori,
\emph{Birational geometry of algebraic varieties},
Cambridge Tracts in Mathematics, vol.~134,
Cambridge University Press, Cambridge, 1998.
\href{https://doi.org/10.1017/CBO9780511662560}{DOI}.

\bibitem{Langer}
A.~Langer,
\emph{Semistable sheaves in positive characteristic},
Ann. of Math. (2) \textbf{159} (2004), no.~1, 251--276.
\href{https://doi.org/10.4007/annals.2004.159.251}{DOI}.

\bibitem{LazarsfeldPAGI}
R.~Lazarsfeld,
\emph{Positivity in Algebraic Geometry I: Classical Setting: Line Bundles and Linear Series},
Ergebnisse der Mathematik und ihrer Grenzgebiete, vol.~48,
Springer-Verlag, 2004.
\href{https://www.math.stonybrook.edu/~roblaz/Reprints/Laz.PAG.Chapt1.pdf}{Author-hosted Chapter~1}.

\bibitem{LP}
V.~Lazi\'c and Th.~Peternell,
\emph{On Generalised Abundance, I},
Publ. Res. Inst. Math. Sci. \textbf{56} (2020), no.~2, 353--389.
\href{https://doi.org/10.4171/PRIMS/56-2-3}{DOI}.
Numbering follows arXiv:1808.00438v2, 6 September 2018.

\bibitem{LPII}
Vladimir Lazi\'c and Thomas Peternell,
\emph{On Generalised Abundance, II},
Peking Math. J. \textbf{3} (2020), no.~1, 1--46.
\href{https://doi.org/10.1007/s42543-019-00022-1}{DOI}.

\bibitem{MR}
V.~B.~Mehta and A.~Ramanathan,
\emph{Semistable sheaves on projective varieties and their restriction to curves},
Math.~Ann. \textbf{258} (1982), no.~3, 213--224.
\href{https://doi.org/10.1007/BF01450677}{DOI}.

\bibitem{Miyaoka}
Y.~Miyaoka,
\emph{Deformations of a morphism along a foliation and applications},
in \emph{Algebraic Geometry, Bowdoin, 1985} (Brunswick, Maine, 1985),
Proc.~Sympos.~Pure Math. \textbf{46}, Part~1,
American Mathematical Society, Providence, RI, 1987, pp.~245--268.
\href{https://doi.org/10.1090/pspum/046.1/927960}{DOI}.

\bibitem{MustataSchwede}
M.~Musta\c{t}\u{a} and K.~Schwede,
\emph{A Frobenius variant of Seshadri constants},
Math. Ann. \textbf{358} (2014), nos.~3--4, 861--878.
\href{https://doi.org/10.1007/s00208-013-0976-4}{DOI}.
Numbering follows arXiv:1203.1081v1.

\bibitem{Nakayama}
N.~Nakayama,
\emph{Zariski-decomposition and abundance},
MSJ Memoirs, vol.~14,
Mathematical Society of Japan, Tokyo, 2004.

\bibitem{LA}
OpenAI,
\emph{Log abundance in characteristic zero},
OpenAI Math Release preprint
\href{https://github.com/openai/math/blob/main/preprints/Log-abundance-in-characteristic-zero-September-24-2026/paper.pdf}{OAI:Log-abundance-in-characteristic-zero-September-24-2026},
2026.

\bibitem{MM}
OpenAI,
\emph{Minimal models and Mori fibre spaces for generalized log canonical $\mathbb{Q}$-pairs},
OpenAI Math Release preprint
\href{https://github.com/openai/math/blob/main/preprints/Minimal-models-and-Mori-fibre-spaces-for-generalized-log-canonical-Q-pairs-September-24-2026/paper.pdf}{OAI:Minimal-models-and-Mori-fibre-spaces-for-generalized-log-canonical-Q-pairs-September-24-2026},
2026.

\bibitem{Sun}
X.~Sun,
\emph{Direct images of bundles under Frobenius morphism},
Invent. Math. \textbf{173} (2008), no.~2, 427--447.
\href{https://doi.org/10.1007/s00222-008-0125-y}{DOI}.
Numbering follows arXiv:math/0611360v2.

\end{thebibliography}
\end{document}
